\hypertarget{csl-1.0.2-specification}{%
\section{CSL 1.0.2 Specification}\label{csl-1.0.2-specification}}

Principal Authors \href{https://twitter.com/rintzezelle}{Rintze M.
Zelle, PhD}, \href{https://twitter.com/bmwiernik}{Brenton M. Wiernik},
Frank G. Bennett, Jr., Bruce D\textquotesingle Arcus, Denis Maier

with additional contributions from Julien Gonzalez, Sebastian Karcher,
Sylvester Keil, Cormac Relf, Lars Willighagen, and other CSL
contributors.

\includegraphics{/media/cc-by-sa-80x15.png}\_\\
This work is licensed under a
\href{https://creativecommons.org/licenses/by-sa/4.0/}{Creative Commons
Attribution-ShareAlike 4.0 International License}.

\textbf{Table of Contents}

\hypertarget{introduction}{%
\subsection{Introduction}\label{introduction}}

The Citation Style Language (CSL) is an XML-based format to describe the
formatting of citations, notes and bibliographies, offering:

\begin{itemize}
\tightlist
\item
  An open format
\item
  Compact and robust styles
\item
  Extensive support for style requirements
\item
  Automatic style localization
\item
  Infrastructure for style distribution and updating
\item
  Thousands of freely available styles (Creative Commons BY-SA licensed)
\end{itemize}

For additional documentation, the CSL schema, styles, and locales, visit
the CSL project home,
\href{http://citationstyles.org}{citationstyles.org}.

\hypertarget{terminology}{%
\subsubsection{Terminology}\label{terminology}}

The key words MUST, MUST NOT, REQUIRED, SHALL, SHALL NOT, SHOULD, SHOULD
NOT, RECOMMENDED, MAY, and OPTIONAL, are to be interpreted as described
in \href{http://tools.ietf.org/html/rfc2119}{IETF RFC 2119}.

\hypertarget{namespacing}{%
\subsection{Namespacing}\label{namespacing}}

The CSL \href{http://en.wikipedia.org/wiki/XML_Namespace}{XML namespace
URI} is "\url{http://purl.org/net/xbiblio/csl}". The namespace prefix
\texttt{cs:} is used throughout this specification when referring to CSL
elements, but is generally omitted in favor of a default namespace
declaration (set with the \texttt{xmlns} attribute) on the root
\texttt{cs:style} or \texttt{cs:locale} element.

\hypertarget{whitespace-handling}{%
\subsection{Whitespace Handling}\label{whitespace-handling}}

CSL styles are valid XML, but CSL processors MUST NOT normalize
attribute values by trimming leading or trailing whitespace from
attributes which define text that is intended for output:

\begin{itemize}
\tightlist
\item
  after-collapse-delimiter
\item
  cite-group-delimiter
\item
  delimiter
\item
  initialize-with
\item
  name-delimiter
\item
  names-delimiter
\item
  prefix
\item
  range-delimiter
\item
  sort-separator
\item
  suffix
\item
  year-suffix-delimiter
\item
  value
\end{itemize}

\hypertarget{file-types}{%
\subsection{File Types}\label{file-types}}

There are three types of CSL files: independent and dependent styles
(both types use the ".csl" extension), and locale files (named
"locales-xx-XX.xml", where "xx-XX" is a language dialect, e.g. "en-US"
for American English).

\hypertarget{independent-styles}{%
\subsubsection{Independent Styles}\label{independent-styles}}

Independent styles contain formatting instructions for citations, notes
and bibliographies. While mostly self-contained, they rely on locale
files for (default) localization data.

\hypertarget{dependent-styles}{%
\subsubsection{Dependent Styles}\label{dependent-styles}}

A dependent style is an alias for an independent style. Its contents are
limited to style metadata, and doesn\textquotesingle t include any
formatting instructions (the sole exception is that dependent styles can
specify an overriding style locale). By linking dependent styles for
journals that share the same citation style (e.g., "Nature
Biotechnology", "Nature Nanotechnology", etc.) to a single independent
style (e.g., "Nature Journals"), there is no need to duplicate
formatting instructions.

\hypertarget{locale-files}{%
\subsubsection{Locale Files}\label{locale-files}}

Each locale file contains a set of localization data (term translations,
localized date formats, and grammar options) for a particular language
dialect.

\hypertarget{xml-declaration}{%
\subsection{XML Declaration}\label{xml-declaration}}

Each style or locale should begin with an XML declaration, specifying
the XML version and character encoding. In most cases, the declaration
will be:

\begin{Shaded}
\begin{Highlighting}[]
\FunctionTok{\textless{}?xml}\OtherTok{ version=}\StringTok{"1.0"}\OtherTok{ encoding=}\StringTok{"UTF{-}8"}\FunctionTok{?\textgreater{}}
\end{Highlighting}
\end{Shaded}

\hypertarget{styles---structure}{%
\subsection{Styles - Structure}\label{styles---structure}}

\hypertarget{the-root-element---csstyle}{%
\subsubsection{\texorpdfstring{The Root Element -
\texttt{cs:style}}{The Root Element - cs:style}}\label{the-root-element---csstyle}}

The root element of styles is \texttt{cs:style}. In independent styles,
the element carries the following attributes:

\begin{description}
\item[\texttt{class}]
Determines whether the style uses in-text citations (value "in-text") or
notes ("note").
\item[\texttt{default-locale} (optional)]
Sets a default locale for style localization. Value must be a
\href{http://books.xmlschemata.org/relaxng/ch19-77191.html}{locale
code}.
\item[\texttt{version}]
The CSL version of the style. Must be "1.0" for CSL 1.0-compatible
styles.
\end{description}

In addition, \texttt{cs:style} may carry any of the
\protect\hyperlink{global-options}{global options} and
\protect\hyperlink{inheritable-name-options}{inheritable name options}.

Of these attributes, only \texttt{version} is required on
\texttt{cs:style} in dependent styles, while the \texttt{default-locale}
attribute may be set to specify an overriding style locale. The other
attributes are allowed but ignored.

An example of \texttt{cs:style} for an independent style, preceded by
the XML declaration:

\begin{Shaded}
\begin{Highlighting}[]
\FunctionTok{\textless{}?xml}\OtherTok{ version=}\StringTok{"1.0"}\OtherTok{ encoding=}\StringTok{"UTF{-}8"}\FunctionTok{?\textgreater{}}
\NormalTok{\textless{}}\KeywordTok{style}\OtherTok{ xmlns=}\StringTok{"http://purl.org/net/xbiblio/csl"}\OtherTok{ version=}\StringTok{"1.0"}\OtherTok{ class=}\StringTok{"in{-}text"}\OtherTok{ default{-}locale=}\StringTok{"fr{-}FR"}\NormalTok{/\textgreater{}}
\end{Highlighting}
\end{Shaded}

\hypertarget{child-elements-of-csstyle}{%
\subsubsection{\texorpdfstring{Child Elements of
\texttt{cs:style}}{Child Elements of cs:style}}\label{child-elements-of-csstyle}}

In independent styles, the \texttt{cs:style} root element has the
following child elements:

\begin{description}
\item[\texttt{cs:info}]
Must appear as the first child element of \texttt{cs:style}. Contains
the metadata describing the style (style name, ID, authors, etc.).
\item[\texttt{cs:citation}]
Must appear once. Describes the formatting of in-text citations or
notes.
\item[\texttt{cs:bibliography} (optional)]
May appear once. Describes the formatting of the bibliography.
\item[\texttt{cs:macro} (optional)]
May appear multiple times. Macros allow formatting instructions to be
reused, keeping styles compact and maintainable.
\item[\texttt{cs:locale} (optional)]
May appear multiple times. Used to specify (overriding) localization
data.
\end{description}

In \protect\hyperlink{dependent-styles}{dependent styles},
\texttt{cs:style} has only one child element, \texttt{cs:info}.

\hypertarget{info}{%
\paragraph{Info}\label{info}}

The \texttt{cs:info} element contains the style\textquotesingle s
metadata. Its structure is based on the
\href{http://tools.ietf.org/html/rfc4287}{Atom Syndication Format}.

In independent styles, \texttt{cs:info} has the following child
elements:

\begin{description}
\item[\texttt{cs:author} and \texttt{cs:contributor} (optional)]
\texttt{cs:author} and \texttt{cs:contributor}, used to respectively
acknowledge style authors and contributors, may each be used multiple
times. Within these elements, the child element \texttt{cs:name} must
appear once, while \texttt{cs:email} and \texttt{cs:uri} each may appear
once. These child elements should contain respectively the name, email
address and URI of the author or contributor.
\item[\texttt{cs:category} (optional)]
Styles may be assigned one or more categories. \texttt{cs:category} may
be used once to describe how in-text citations are rendered, using the
\texttt{citation-format} attribute set to one of the following values:

\begin{itemize}
\tightlist
\item
  "author-date" - e.g. "\ldots{} (Doe, 1999)"
\item
  "author" - e.g. "\ldots{} (Doe)"
\item
  "numeric" - e.g. "\ldots{} {[}1{]}"
\item
  "label" - e.g. "\ldots{} {[}doe99{]}"
\item
  "note" - the citation appears as a footnote or endnote
\end{itemize}

\texttt{cs:category} may be used multiple times with the \texttt{field}
attribute, set to one of the discipline categories (see
\protect\hyperlink{appendix-i---categories}{Appendix I - Categories}),
to indicates the field(s) for which the style is relevant.
\item[\texttt{cs:id}]
Must appear once and contain a stable, unique identifier to establish
the identity of the style. For historical reasons, existing styles may
use URIs, but new styles should use a UUID to guarantee stability and
uniqueness.
\item[\texttt{cs:issn}/\texttt{cs:eissn}/\texttt{cs:issnl} (optional)]
The \texttt{cs:issn} element may be used multiple times to indicate the
ISSN identifier(s) of the journal for which the style was written. The
\texttt{cs:eissn} and \texttt{cs:issnl} elements may each be used once
for the eISSN and
\href{http://www.issn.org/2-22637-What-is-an-ISSN-L.php}{ISSN-L}
identifiers, respectively.
\item[\texttt{cs:link} (optional)]
May be used multiple times. \texttt{cs:link} must carry two attributes:
\texttt{href}, set to a URI (usually a URL), and \texttt{rel}, whose
value indicates how the URI relates to the style. The possible values of
\texttt{rel}:

\begin{itemize}
\tightlist
\item
  "self" - style URI
\item
  "template" - URI of the style from which the current style is derived
\item
  "documentation" - URI of style documentation
\end{itemize}

The \texttt{cs:link} element may contain content describing the link.
\item[\texttt{cs:published} (optional)]
May appear once. The contents of \texttt{cs:published} must be a
\href{http://books.xmlschemata.org/relaxng/ch19-77049.html}{timestamp},
indicating when the style was initially created or made available.
\item[\texttt{cs:rights} (optional)]
May appear once. The contents of \texttt{cs:rights} specifies the
license under which the style file is released. The element may carry a
\texttt{license} attribute to specify the URI of the license.
\item[\texttt{cs:summary} (optional)]
May appear once. The contents of \texttt{cs:summary} gives a (short)
description of the style.
\item[\texttt{cs:title}]
Must appear once. The contents of \texttt{cs:title} should be the name
of the style as shown to users.
\item[\texttt{cs:title-short} (optional)]
May appear once. The contents of \texttt{cs:title-short} should be a
shortened style name (e.g. "APA").
\item[\texttt{cs:updated}]
Must appear once. The contents of \texttt{cs:updated} must be a
\href{http://books.xmlschemata.org/relaxng/ch19-77049.html}{timestamp}
that shows when the style was last updated.
\end{description}

The \texttt{cs:link}, \texttt{cs:rights}, \texttt{cs:summary},
\texttt{cs:title} and \texttt{cs:title-short} elements may carry a
\texttt{xml:lang} attribute to specify the language of the
element\textquotesingle s content (the value must be an
\href{http://books.xmlschemata.org/relaxng/ch19-77191.html}{xsd:language
locale code}). For \texttt{cs:link}, the attribute can also be used to
indicate the language of the link target.

In \protect\hyperlink{dependent-styles}{dependent styles},
\texttt{cs:link} must be used with \texttt{rel} set to
"independent-parent", with the URI of the independent parent style set
on \texttt{href}. In addition, \texttt{cs:link} may not be used with
\texttt{rel} set to "template".

An example of \texttt{cs:info} for an independent style:

\begin{Shaded}
\begin{Highlighting}[]
\NormalTok{\textless{}}\KeywordTok{info}\NormalTok{\textgreater{}}
\NormalTok{  \textless{}}\KeywordTok{title}\NormalTok{\textgreater{}Style Title\textless{}/}\KeywordTok{title}\NormalTok{\textgreater{}}
\NormalTok{  \textless{}}\KeywordTok{id}\NormalTok{\textgreater{}http://www.zotero.org/styles/style{-}title\textless{}/}\KeywordTok{id}\NormalTok{\textgreater{}}
\NormalTok{  \textless{}}\KeywordTok{link}\OtherTok{ href=}\StringTok{"http://www.zotero.org/styles/style{-}title"}\OtherTok{ rel=}\StringTok{"self"}\NormalTok{/\textgreater{}}
\NormalTok{  \textless{}}\KeywordTok{link}\OtherTok{ href=}\StringTok{"http://www.example.org/instructions{-}to{-}authors\#references"}\OtherTok{ rel=}\StringTok{"documentation"}\NormalTok{/\textgreater{}}
\NormalTok{  \textless{}}\KeywordTok{author}\NormalTok{\textgreater{}}
\NormalTok{    \textless{}}\KeywordTok{name}\NormalTok{\textgreater{}Author Name\textless{}/}\KeywordTok{name}\NormalTok{\textgreater{}}
\NormalTok{    \textless{}}\KeywordTok{email}\NormalTok{\textgreater{}name@example.org\textless{}/}\KeywordTok{email}\NormalTok{\textgreater{}}
\NormalTok{    \textless{}}\KeywordTok{uri}\NormalTok{\textgreater{}http://www.example.org/name\textless{}/}\KeywordTok{uri}\NormalTok{\textgreater{}}
\NormalTok{  \textless{}/}\KeywordTok{author}\NormalTok{\textgreater{}}
\NormalTok{  \textless{}}\KeywordTok{category}\OtherTok{ citation{-}format=}\StringTok{"author{-}date"}\NormalTok{/\textgreater{}}
\NormalTok{  \textless{}}\KeywordTok{category}\OtherTok{ field=}\StringTok{"zoology"}\NormalTok{/\textgreater{}}
\NormalTok{  \textless{}}\KeywordTok{updated}\NormalTok{\textgreater{}2011{-}10{-}29T21:01:24+00:00\textless{}/}\KeywordTok{updated}\NormalTok{\textgreater{}}
\NormalTok{  \textless{}}\KeywordTok{rights}\OtherTok{ license=}\StringTok{"http://creativecommons.org/licenses/by{-}sa/3.0/"}\NormalTok{\textgreater{}This work}
\NormalTok{  is licensed under a Creative Commons Attribution{-}ShareAlike 3.0 License\textless{}/}\KeywordTok{rights}\NormalTok{\textgreater{}}
\NormalTok{\textless{}/}\KeywordTok{info}\NormalTok{\textgreater{}}
\end{Highlighting}
\end{Shaded}

\hypertarget{citation}{%
\paragraph{Citation}\label{citation}}

The \texttt{cs:citation} element describes the formatting of citations,
which consist of one or more references ("cites") to bibliographic
sources. Citations appear in the form of either in-text citations (in
the author (e.g. "{[}Doe{]}"), author-date ("{[}Doe 1999{]}"), label
("{[}doe99{]}") or number ("{[}1{]}") format) or notes. The required
\texttt{cs:layout} child element describes what, and how, bibliographic
data should be included in the citations (see
\protect\hyperlink{layout}{Layout}). \texttt{cs:layout} may be preceded
by a \texttt{cs:sort} element, which can be used to specify how cites
within a citation should be sorted (see
\protect\hyperlink{sorting}{Sorting}). The \texttt{cs:citation} element
may carry attributes for
\protect\hyperlink{citation-specific-options}{Citation-specific Options}
and \protect\hyperlink{inheritable-name-options}{Inheritable Name
Options}. An example of a \texttt{cs:citation} element:

\begin{Shaded}
\begin{Highlighting}[]
\NormalTok{\textless{}}\KeywordTok{citation}\NormalTok{\textgreater{}}
\NormalTok{  \textless{}}\KeywordTok{sort}\NormalTok{\textgreater{}}
\NormalTok{    \textless{}}\KeywordTok{key}\OtherTok{ variable=}\StringTok{"citation{-}number"}\NormalTok{/\textgreater{}}
\NormalTok{  \textless{}/}\KeywordTok{sort}\NormalTok{\textgreater{}}
\NormalTok{  \textless{}}\KeywordTok{layout}\NormalTok{\textgreater{}}
\NormalTok{    \textless{}}\KeywordTok{text}\OtherTok{ variable=}\StringTok{"citation{-}number"}\NormalTok{/\textgreater{}}
\NormalTok{  \textless{}/}\KeywordTok{layout}\NormalTok{\textgreater{}}
\NormalTok{\textless{}/}\KeywordTok{citation}\NormalTok{\textgreater{}}
\end{Highlighting}
\end{Shaded}

\textbf{A note to CSL processor developers} In note styles, a citation
is often a sentence by itself. Therefore, the first character of a
citation should preferably be uppercased when there is no preceding text
in the note. In all other cases (e.g. when a citation is inserted into
the middle of a pre-existing footnote), the citation should be printed
as is.

\hypertarget{bibliography}{%
\paragraph{Bibliography}\label{bibliography}}

The \texttt{cs:bibliography} element describes the formatting of
bibliographies, which list one or more bibliographic sources. The
required \texttt{cs:layout} child element describes how each
bibliographic entry should be formatted. \texttt{cs:layout} may be
preceded by a \texttt{cs:sort} element, which can be used to specify how
references within the bibliography should be sorted (see
\protect\hyperlink{sorting}{Sorting}). The \texttt{cs:bibliography}
element may carry attributes for
\protect\hyperlink{bibliography-specific-options}{Bibliography-specific
Options} and \protect\hyperlink{inheritable-name-options}{Inheritable
Name Options}. An example of a \texttt{cs:bibliography} element:

\begin{Shaded}
\begin{Highlighting}[]
\NormalTok{\textless{}}\KeywordTok{bibliography}\NormalTok{\textgreater{}}
\NormalTok{  \textless{}}\KeywordTok{sort}\NormalTok{\textgreater{}}
\NormalTok{    \textless{}}\KeywordTok{key}\OtherTok{ macro=}\StringTok{"author"}\NormalTok{/\textgreater{}}
\NormalTok{  \textless{}/}\KeywordTok{sort}\NormalTok{\textgreater{}}
\NormalTok{  \textless{}}\KeywordTok{layout}\NormalTok{\textgreater{}}
\NormalTok{    \textless{}}\KeywordTok{group}\OtherTok{ delimiter=}\StringTok{". "}\NormalTok{\textgreater{}}
\NormalTok{      \textless{}}\KeywordTok{text}\OtherTok{ macro=}\StringTok{"author"}\NormalTok{/\textgreater{}}
\NormalTok{      \textless{}}\KeywordTok{text}\OtherTok{ variable=}\StringTok{"title"}\NormalTok{/\textgreater{}}
\NormalTok{    \textless{}/}\KeywordTok{group}\NormalTok{\textgreater{}}
\NormalTok{  \textless{}/}\KeywordTok{layout}\NormalTok{\textgreater{}}
\NormalTok{\textless{}/}\KeywordTok{bibliography}\NormalTok{\textgreater{}}
\end{Highlighting}
\end{Shaded}

\hypertarget{macro}{%
\paragraph{Macro}\label{macro}}

Macros, defined with \texttt{cs:macro} elements, contain formatting
instructions. Macros can be called with \texttt{cs:text} from within
other macros and the \texttt{cs:layout} element of \texttt{cs:citation}
and \texttt{cs:bibliography}, and with \texttt{cs:key} from within
\texttt{cs:sort} of \texttt{cs:citation} and \texttt{cs:bibliography}.
It is recommended to place macros after any \texttt{cs:locale} elements
and before the \texttt{cs:citation} element.

Macros are referenced by the value of the required \texttt{name}
attribute on \texttt{cs:macro}. The \texttt{cs:macro} element must
contain one or more \protect\hyperlink{rendering-elements}{rendering
elements}.

The use of macros can improve style readability, compactness and
maintainability. It is recommended to keep the contents of
\texttt{cs:citation} and \texttt{cs:bibliography} compact and agnostic
of item types (e.g. books, journal articles, etc.) by depending on macro
calls. To allow for easy reuse of macros in other styles, it is
recommended to use common macro names.

In the example below, cites consist of the item title, rendered in
italics when the item type is "book":

\begin{Shaded}
\begin{Highlighting}[]
\NormalTok{\textless{}}\KeywordTok{style}\NormalTok{\textgreater{}}
\NormalTok{  \textless{}}\KeywordTok{macro}\OtherTok{ name=}\StringTok{"title"}\NormalTok{\textgreater{}}
\NormalTok{    \textless{}}\KeywordTok{choose}\NormalTok{\textgreater{}}
\NormalTok{      \textless{}}\KeywordTok{if}\OtherTok{ type=}\StringTok{"book"}\NormalTok{\textgreater{}}
\NormalTok{        \textless{}}\KeywordTok{text}\OtherTok{ variable=}\StringTok{"title"}\OtherTok{ font{-}style=}\StringTok{"italic"}\NormalTok{/\textgreater{}}
\NormalTok{      \textless{}/}\KeywordTok{if}\NormalTok{\textgreater{}}
\NormalTok{      \textless{}}\KeywordTok{else}\NormalTok{\textgreater{}}
\NormalTok{        \textless{}}\KeywordTok{text}\OtherTok{ variable=}\StringTok{"title"}\NormalTok{/\textgreater{}}
\NormalTok{      \textless{}/}\KeywordTok{else}\NormalTok{\textgreater{}}
\NormalTok{    \textless{}/}\KeywordTok{choose}\NormalTok{\textgreater{}}
\NormalTok{  \textless{}/}\KeywordTok{macro}\NormalTok{\textgreater{}}
\NormalTok{  \textless{}}\KeywordTok{citation}\NormalTok{\textgreater{}}
\NormalTok{    \textless{}}\KeywordTok{layout}\NormalTok{\textgreater{}}
\NormalTok{      \textless{}}\KeywordTok{text}\OtherTok{ macro=}\StringTok{"title"}\NormalTok{/\textgreater{}}
\NormalTok{    \textless{}/}\KeywordTok{layout}\NormalTok{\textgreater{}}
\NormalTok{  \textless{}/}\KeywordTok{citation}\NormalTok{\textgreater{}}
\NormalTok{\textless{}/}\KeywordTok{style}\NormalTok{\textgreater{}}
\end{Highlighting}
\end{Shaded}

Delimiters from any ancestor delimiting element are not applied within
the output of a \texttt{\textless{}text\ macro="..."\textgreater{}}
element (see \protect\hyperlink{delimiter}{delimiter}).

\hypertarget{locale}{%
\paragraph{Locale}\label{locale}}

Localization data, by default drawn from the "locales-xx-XX.xml" locale
files, may be redefined or supplemented with \texttt{cs:locale}
elements, which should be placed directly after the \texttt{cs:info}
element.

The value of the optional \texttt{xml:lang} attribute on
\texttt{cs:locale}, which must be set to an
\href{http://books.xmlschemata.org/relaxng/ch19-77191.html}{xsd:language
locale code}, determines which languages or language dialects are
affected (see \protect\hyperlink{locale-fallback}{Locale Fallback}).

See \protect\hyperlink{terms}{Terms},
\protect\hyperlink{localized-date-formats}{Localized Date Formats} and
\protect\hyperlink{localized-options}{Localized Options} for further
details on the use of \texttt{cs:locale}.

An example of \texttt{cs:locale} in a style:

\begin{Shaded}
\begin{Highlighting}[]
\NormalTok{\textless{}}\KeywordTok{style}\NormalTok{\textgreater{}}
\NormalTok{  \textless{}}\KeywordTok{locale}\OtherTok{ xml:lang=}\StringTok{"en"}\NormalTok{\textgreater{}}
\NormalTok{    \textless{}}\KeywordTok{terms}\NormalTok{\textgreater{}}
\NormalTok{      \textless{}}\KeywordTok{term}\OtherTok{ name=}\StringTok{"editortranslator"}\OtherTok{ form=}\StringTok{"short"}\NormalTok{\textgreater{}}
\NormalTok{        \textless{}}\KeywordTok{single}\NormalTok{\textgreater{}ed. }\DecValTok{\&amp;}\NormalTok{ trans.\textless{}/}\KeywordTok{single}\NormalTok{\textgreater{}}
\NormalTok{        \textless{}}\KeywordTok{multiple}\NormalTok{\textgreater{}eds. }\DecValTok{\&amp;}\NormalTok{ trans.\textless{}/}\KeywordTok{multiple}\NormalTok{\textgreater{}}
\NormalTok{      \textless{}/}\KeywordTok{term}\NormalTok{\textgreater{}}
\NormalTok{    \textless{}/}\KeywordTok{terms}\NormalTok{\textgreater{}}
\NormalTok{  \textless{}/}\KeywordTok{locale}\NormalTok{\textgreater{}}
\NormalTok{\textless{}/}\KeywordTok{style}\NormalTok{\textgreater{}}
\end{Highlighting}
\end{Shaded}

\hypertarget{locale-fallback}{%
\subparagraph{Locale Fallback}\label{locale-fallback}}

Locale files provide localization data for language dialects (e.g.
"en-US" for American English), whereas the optional \texttt{cs:locale}
elements in styles can either lack the \texttt{xml:lang} attribute, or
have it set to either a language (e.g. "en" for English) or dialect.
Locale fallback is the mechanism determining from which of these sources
each localizable unit (a date format, localized option, or specific form
of a term) is retrieved.

For dialects of the same language, one is designated the primary
dialect. All others are secondaries. At the moment of writing, the
available locale files include:

\begin{longtable}[]{@{}
  >{\raggedright\arraybackslash}p{(\columnwidth - 2\tabcolsep) * \real{0.2917}}
  >{\raggedright\arraybackslash}p{(\columnwidth - 2\tabcolsep) * \real{0.5417}}@{}}
\toprule\noalign{}
\begin{minipage}[b]{\linewidth}\raggedright
Primary dialect
\end{minipage} & \begin{minipage}[b]{\linewidth}\raggedright
Secondary dialect(s)
\end{minipage} \\
\midrule\noalign{}
\endhead
\bottomrule\noalign{}
\endlastfoot
de-DE (German) & de-AT (Austria), de-CH (Switzerland) \\
en-US (English) & en-GB (UK) \\
es-ES (Spanish) & es-CL (Chile), es-MX (Mexico) \\
fr-FR (French) & fr-CA (Canada) \\
pt-PT (Portuguese) & pt-BR (Brazil) \\
zh-CN (Chinese) & zh-TW (Taiwan) \\
\end{longtable}

Locale fallback is best described with an example. If the chosen output
locale is "de-AT" (Austrian German), localizable units are individually
drawn from the following sources, in decreasing order of priority:

\begin{enumerate}
\def\labelenumi{\Alph{enumi}.}
\tightlist
\item
  In-style \texttt{cs:locale} elements

  \begin{enumerate}
  \def\labelenumii{\arabic{enumii}.}
  \tightlist
  \item
    \texttt{xml:lang} set to chosen dialect, "de-AT"
  \item
    \texttt{xml:lang} set to matching language, "de" (German)
  \item
    \texttt{xml:lang} not set
  \end{enumerate}
\item
  Locale files

  \begin{enumerate}
  \def\labelenumii{\arabic{enumii}.}
  \setcounter{enumii}{3}
  \tightlist
  \item
    \texttt{xml:lang} set to chosen dialect, "de-AT"
  \item
    \texttt{xml:lang} set to matching primary dialect, "de-DE" (Standard
    German) (only applicable when the chosen locale is a secondary
    dialect)
  \item
    \texttt{xml:lang} set to "en-US" (American English)
  \end{enumerate}
\end{enumerate}

If the chosen output locale is a language (e.g. "de"), the (primary)
dialect is used in step 1 (e.g. "de-DE").

Fallback stops once a localizable unit has been found. For terms, this
even is the case when they are defined as empty strings (e.g.
\texttt{\textless{}term\ name="and"/\textgreater{}} or
\texttt{\textless{}term\ name="and"\textgreater{}\textless{}/term\textgreater{}}).
Locale fallback takes precedence over fallback of term forms (see
\protect\hyperlink{terms}{Terms}).

\hypertarget{locale-files---structure}{%
\subsection{Locale Files - Structure}\label{locale-files---structure}}

While localization data can be included in styles (see
\protect\hyperlink{locale}{Locale}), locale files conveniently provide
sets of default localization data, consisting of terms, date formats and
grammar options.

Each locale file contains localization data for a single language
dialect. This
\href{http://books.xmlschemata.org/relaxng/ch19-77191.html}{locale code}
is set on the required \texttt{xml:lang} attribute on the
\texttt{cs:locale} root element. The same locale code must also be used
in the file name of the locale file (the "xx-XX" in
"locales-xx-XX.xml"). The root element must carry the \texttt{version}
attribute, indicating the CSL version of the locale file (must be "1.0"
for CSL 1.0-compatible locale files). Locale files have the same
requirements for \protect\hyperlink{namespacing}{namespacing} as styles.
The \texttt{cs:locale} element may contain \texttt{cs:info} as its first
child element, and requires the child elements \texttt{cs:terms},
\texttt{cs:date} and \texttt{cs:style-options} (these elements are
described below). An example showing part of a locale file:

\begin{Shaded}
\begin{Highlighting}[]
\FunctionTok{\textless{}?xml}\OtherTok{ version=}\StringTok{"1.0"}\OtherTok{ encoding=}\StringTok{"UTF{-}8"}\FunctionTok{?\textgreater{}}
\NormalTok{\textless{}}\KeywordTok{locale}\OtherTok{ xml:lang=}\StringTok{"en{-}US"}\OtherTok{ version=}\StringTok{"1.0"}\OtherTok{ xmlns=}\StringTok{"http://purl.org/net/xbiblio/csl"}\NormalTok{\textgreater{}}
\NormalTok{  \textless{}}\KeywordTok{style{-}options}\OtherTok{ punctuation{-}in{-}quote=}\StringTok{"true"}\NormalTok{/\textgreater{}}
\NormalTok{  \textless{}}\KeywordTok{date}\OtherTok{ form=}\StringTok{"text"}\NormalTok{\textgreater{}}
\NormalTok{    \textless{}}\KeywordTok{date{-}part}\OtherTok{ name=}\StringTok{"month"}\OtherTok{ suffix=}\StringTok{" "}\NormalTok{/\textgreater{}}
\NormalTok{    \textless{}}\KeywordTok{date{-}part}\OtherTok{ name=}\StringTok{"day"}\OtherTok{ suffix=}\StringTok{", "}\NormalTok{/\textgreater{}}
\NormalTok{    \textless{}}\KeywordTok{date{-}part}\OtherTok{ name=}\StringTok{"year"}\NormalTok{/\textgreater{}}
\NormalTok{  \textless{}/}\KeywordTok{date}\NormalTok{\textgreater{}}
\NormalTok{  \textless{}}\KeywordTok{date}\OtherTok{ form=}\StringTok{"numeric"}\NormalTok{\textgreater{}}
\NormalTok{    \textless{}}\KeywordTok{date{-}part}\OtherTok{ name=}\StringTok{"year"}\NormalTok{/\textgreater{}}
\NormalTok{    \textless{}}\KeywordTok{date{-}part}\OtherTok{ name=}\StringTok{"month"}\OtherTok{ form=}\StringTok{"numeric"}\OtherTok{ prefix=}\StringTok{"{-}"}\OtherTok{ range{-}delimiter=}\StringTok{"/"}\NormalTok{/\textgreater{}}
\NormalTok{    \textless{}}\KeywordTok{date{-}part}\OtherTok{ name=}\StringTok{"day"}\OtherTok{ prefix=}\StringTok{"{-}"}\OtherTok{ range{-}delimiter=}\StringTok{"/"}\NormalTok{/\textgreater{}}
\NormalTok{  \textless{}/}\KeywordTok{date}\NormalTok{\textgreater{}}
\NormalTok{  \textless{}}\KeywordTok{terms}\NormalTok{\textgreater{}}
\NormalTok{    \textless{}}\KeywordTok{term}\OtherTok{ name=}\StringTok{"no date"}\NormalTok{\textgreater{}n.d.\textless{}/}\KeywordTok{term}\NormalTok{\textgreater{}}
\NormalTok{    \textless{}}\KeywordTok{term}\OtherTok{ name=}\StringTok{"et{-}al"}\NormalTok{\textgreater{}et al.\textless{}/}\KeywordTok{term}\NormalTok{\textgreater{}}
\NormalTok{    \textless{}}\KeywordTok{term}\OtherTok{ name=}\StringTok{"page"}\NormalTok{\textgreater{}}
\NormalTok{      \textless{}}\KeywordTok{single}\NormalTok{\textgreater{}page\textless{}/}\KeywordTok{single}\NormalTok{\textgreater{}}
\NormalTok{      \textless{}}\KeywordTok{multiple}\NormalTok{\textgreater{}pages\textless{}/}\KeywordTok{multiple}\NormalTok{\textgreater{}}
\NormalTok{    \textless{}/}\KeywordTok{term}\NormalTok{\textgreater{}}
\NormalTok{    \textless{}}\KeywordTok{term}\OtherTok{ name=}\StringTok{"page"}\OtherTok{ form=}\StringTok{"short"}\NormalTok{\textgreater{}}
\NormalTok{      \textless{}}\KeywordTok{single}\NormalTok{\textgreater{}p.\textless{}/}\KeywordTok{single}\NormalTok{\textgreater{}}
\NormalTok{      \textless{}}\KeywordTok{multiple}\NormalTok{\textgreater{}pp.\textless{}/}\KeywordTok{multiple}\NormalTok{\textgreater{}}
\NormalTok{    \textless{}/}\KeywordTok{term}\NormalTok{\textgreater{}}
\NormalTok{  \textless{}/}\KeywordTok{terms}\NormalTok{\textgreater{}}
\NormalTok{\textless{}/}\KeywordTok{locale}\NormalTok{\textgreater{}}
\end{Highlighting}
\end{Shaded}

\hypertarget{info-1}{%
\subsubsection{Info}\label{info-1}}

The \texttt{cs:info} element may be used to give metadata on the locale
file. It has the following child elements:

\begin{description}
\item[\texttt{cs:translator} (optional)]
\texttt{cs:translator}, used to acknowledge locale translators, may be
used multiple times. Within the element, the child element
\texttt{cs:name} must appear once, while \texttt{cs:email} and
\texttt{cs:uri} each may appear once. These child elements should
contain respectively the name, email address and URI of the translator.
\item[\texttt{cs:rights} (optional)]
May appear once. The contents of \texttt{cs:rights} specifies the
license under which the locale file is released. The element may carry a
\texttt{license} attribute to specify the URI of the license, and a
\texttt{xml:lang} attribute to specify the language of the
element\textquotesingle s content (the value must be an
\href{http://books.xmlschemata.org/relaxng/ch19-77191.html}{xsd:language
locale code}).
\item[\texttt{cs:updated} (optional)]
May appear once. The contents of \texttt{cs:updated} must be a
\href{http://books.xmlschemata.org/relaxng/ch19-77049.html}{timestamp}
that shows when the locale file was last updated.
\end{description}

\hypertarget{terms}{%
\subsubsection{Terms}\label{terms}}

Terms are localized strings (e.g. by using the "and" term, "Doe and
Smith" automatically becomes "Doe und Smith" when the style locale is
switched from English to German). Terms are defined with
\texttt{cs:term} elements, child elements of \texttt{cs:terms}. Each
\texttt{cs:term} element must carry a \texttt{name} attribute, set to
one of the terms listed in
\protect\hyperlink{appendix-ii---terms}{Appendix II - Terms}.

Terms are either directly defined in the content of \texttt{cs:term},
or, in cases where singular and plural variants are needed (e.g. "page"
and "pages"), in the content of the child elements \texttt{cs:single}
and \texttt{cs:multiple}, respectively.

Terms may be defined for specific forms by using \texttt{cs:term} with
the optional \texttt{form} attribute set to:

\begin{itemize}
\tightlist
\item
  "long" - (default), e.g. "editor" and "editors" for the "editor" term
\item
  "short" - e.g. "ed." and "eds." for the term "editor"
\item
  "verb" - e.g. "edited by" for the term "editor"
\item
  "verb-short" - e.g. "ed." for the term "editor"
\item
  "symbol" - e.g. "§" and "§§" for the term "section"
\end{itemize}

If a style uses a term in a form that is undefined (even after
\protect\hyperlink{locale-fallback}{Locale Fallback}), there is fallback
to other forms: "verb-short" first falls back to "verb", "symbol" first
falls back to "short", and "verb" and "short" both fall back to "long".
In addition, the terms "long-ordinal-01" to "long-ordinal-10" fall back
to the set of ordinal suffix terms. If no locale or form fallback is
available, the term is rendered as an empty string.

The \texttt{match}, \texttt{gender}, and \texttt{gender-form} attributes
can be used on \texttt{cs:term} for the formatting of number variables
rendered as ordinals (e.g. "first", "2nd"). See
\protect\hyperlink{ordinal-suffixes}{Ordinal Suffixes} and
\protect\hyperlink{gender-specific-ordinals}{Gender-specific Ordinals}
below.

Term content should not contain markup such as LaTeX or HTML.
\href{http://unicode.org/reports/tr30/datafiles/SuperscriptFolding.txt}{Superscripted
Unicode characters} can be used for superscripting.

\hypertarget{ordinal-suffixes}{%
\paragraph{Ordinal Suffixes}\label{ordinal-suffixes}}

Number variables can be rendered with \texttt{cs:number} in the
"ordinal" form, e.g. "2nd" (see \protect\hyperlink{number}{Number}). The
ordinal suffixes ("nd" for "2nd") are defined with terms.

The "ordinal" term defines the default ordinal suffix. This default
suffix may be overridden for certain numbers with the following terms:

\begin{itemize}
\tightlist
\item
  "ordinal-00" through "ordinal-09" - by default, a term in this group
  is used when the last digit in the term name matches the last digit of
  the rendered number. E.g. "ordinal-00" would match the numbers "0",
  "10", "20", etc. By setting the optional \texttt{match} attribute to
  "last-two-digits" ("last-digit" is the default), matches are limited
  to numbers where the two last digits agree ("0", "100", "200", etc.).
  When \texttt{match} is set to "whole-number", there is only a match if
  the number is the same as that of the term.
\item
  "ordinal-10" through "ordinal-99" - by default, a term in this group
  is used when the last two digits in the term name match the last two
  digits of the rendered number. When the optional \texttt{match}
  attribute is set to "whole-number" ("last-two-digits" is the default),
  there is only a match if the number is the same as that of the term.
\end{itemize}

When a number has matching terms from both groups (e.g. "13" can match
"ordinal-03" and "ordinal-13"), the term from the "ordinal-10" through
"ordinal-99" group is used.

Ordinal terms work differently in CSL 1.0.1 and later than they did in
CSL 1.0. When neither the style or locale file define the "ordinal"
term, but do define the terms "ordinal-01" through "ordinal-04", the
original CSL 1.0 scheme is used: "ordinal-01" is used for numbers ending
on a 1 (except those ending on 11), "ordinal-02" for those ending on a 2
(except those ending on 12), "ordinal-03" for those ending on a 3
(except those ending on 13) and "ordinal-04" for all other numbers.

The "ordinal" term, and "ordinal-00" through "ordinal-99" terms, behave
differently from other terms when it comes to {Locale Fallback}. Whereas
other terms can be (re)defined individually, (re)defining any of the
ordinal terms through \texttt{cs:locale} replaces all previously defined
ordinal terms.

\hypertarget{long-ordinals}{%
\paragraph{Long Ordinals}\label{long-ordinals}}

Number variables can be rendered with \texttt{cs:number} in the
"long-ordinal" form, e.g. "second" (see
\protect\hyperlink{number}{Number}). The long ordinal terms (e.g.
"second" for "2") are defined with the "long-ordinal-01" through
"long-ordinal-10" terms.\\
Long ordinal forms are available for the numbers 1 through 10. For other
numbers "long-ordinal" falls back to "ordinal".\\
For the numbers 1 through 10 in "long-ordinal" form, the \texttt{match}
attribute is always treated as "whole-number". For other numbers
rendered in "long-ordinal" form, the optional \texttt{match} attribute
follows the behavior described in
\protect\hyperlink{ordinal-suffixes}{Ordinal Suffixes}
("last-two-digits" is the default).

\hypertarget{gender-specific-ordinals}{%
\paragraph{Gender-specific Ordinals}\label{gender-specific-ordinals}}

Some languages use gender-specific ordinals. For example, the English
"1st" and "first" translate in French to "1\textsuperscript{er}" and
"premier" if the target noun is masculine, and "1\textsuperscript{re}"
and "première" if the noun is feminine.

Feminine and masculine variants of the ordinal terms (see
\protect\hyperlink{ordinals}{Ordinals}) may be specified by setting the
\texttt{gender-form} attribute to "feminine" or "masculine" (the term
without \texttt{gender-form} represents the neuter variant). There are
two types of target nouns: a) the terms accompanying the
\protect\hyperlink{number-variables}{number variables}, and b) the month
terms (see \protect\hyperlink{months}{Months}). The gender of these
nouns may be specified on the "long" (default) form of the term using
the \texttt{gender} attribute (set to "feminine" or "masculine"). When a
number variable is rendered with \texttt{cs:number} in the "ordinal" or
"long-ordinal" form, the ordinal term of the same gender is used, with a
fallback to the neuter variant if the feminine or masculine variant is
undefined. When the "day" date-part is rendered in the "ordinal" form,
the ordinal gender is matched against that of the month term.

The example below gives "1re éd." ("1st ed."), "1er janvier" ("January
1st"), and "3e édition" ("3rd edition"):

\begin{Shaded}
\begin{Highlighting}[]
\FunctionTok{\textless{}?xml}\OtherTok{ version=}\StringTok{"1.0"}\OtherTok{ encoding=}\StringTok{"UTF{-}8"}\FunctionTok{?\textgreater{}}
\NormalTok{\textless{}}\KeywordTok{locale}\OtherTok{ xml:lang=}\StringTok{"fr{-}FR"}\NormalTok{\textgreater{}}
\NormalTok{  \textless{}}\KeywordTok{terms}\NormalTok{\textgreater{}}
\NormalTok{    \textless{}}\KeywordTok{term}\OtherTok{ name=}\StringTok{"edition"}\OtherTok{ gender=}\StringTok{"feminine"}\NormalTok{\textgreater{}}
\NormalTok{      \textless{}}\KeywordTok{single}\NormalTok{\textgreater{}édition\textless{}/}\KeywordTok{single}\NormalTok{\textgreater{}}
\NormalTok{      \textless{}}\KeywordTok{multiple}\NormalTok{\textgreater{}éditions\textless{}/}\KeywordTok{multiple}\NormalTok{\textgreater{}}
\NormalTok{    \textless{}/}\KeywordTok{term}\NormalTok{\textgreater{}}
\NormalTok{    \textless{}}\KeywordTok{term}\OtherTok{ name=}\StringTok{"edition"}\OtherTok{ form=}\StringTok{"short"}\NormalTok{\textgreater{}éd.\textless{}/}\KeywordTok{term}\NormalTok{\textgreater{}}
\NormalTok{    \textless{}}\KeywordTok{term}\OtherTok{ name=}\StringTok{"month{-}01"}\OtherTok{ gender=}\StringTok{"masculine"}\NormalTok{\textgreater{}janvier\textless{}/}\KeywordTok{term}\NormalTok{\textgreater{}}
\NormalTok{    \textless{}}\KeywordTok{term}\OtherTok{ name=}\StringTok{"ordinal"}\NormalTok{\textgreater{}e\textless{}/}\KeywordTok{term}\NormalTok{\textgreater{}}
\NormalTok{    \textless{}}\KeywordTok{term}\OtherTok{ name=}\StringTok{"ordinal{-}01"}\OtherTok{ gender{-}form=}\StringTok{"feminine"}\OtherTok{ match=}\StringTok{"whole{-}number"}\NormalTok{\textgreater{}re\textless{}/}\KeywordTok{term}\NormalTok{\textgreater{}}
\NormalTok{    \textless{}}\KeywordTok{term}\OtherTok{ name=}\StringTok{"ordinal{-}01"}\OtherTok{ gender{-}form=}\StringTok{"masculine"}\OtherTok{ match=}\StringTok{"whole{-}number"}\NormalTok{\textgreater{}er\textless{}/}\KeywordTok{term}\NormalTok{\textgreater{}}
\NormalTok{  \textless{}/}\KeywordTok{terms}\NormalTok{\textgreater{}}
\NormalTok{\textless{}/}\KeywordTok{locale}\NormalTok{\textgreater{}}
\end{Highlighting}
\end{Shaded}

\hypertarget{localized-date-formats}{%
\subsubsection{Localized Date Formats}\label{localized-date-formats}}

Two localized date formats can be defined with \texttt{cs:date}
elements: a "numeric" (e.g. "12-15-2005") and a "text" format (e.g.
"December 15, 2005"). The format is set on \texttt{cs:date} with the
required \texttt{form} attribute.

A date format is constructed using \texttt{cs:date-part} child elements
(see \protect\hyperlink{date-part}{Date-part}). With a required
\texttt{name} attribute set to either \texttt{day}, \texttt{month} or
\texttt{year}, the order of these elements reflects the display order of
respectively the day, month, and year. The date can be formatted with
\protect\hyperlink{formatting}{formatting} and
\protect\hyperlink{text-case}{text-case} attributes on the
\texttt{cs:date} and \texttt{cs:date-part} elements. The
\protect\hyperlink{delimiter}{delimiter} attribute may be set on
\texttt{cs:date} to specify the delimiter for the \texttt{cs:date-part}
elements, and \protect\hyperlink{affixes}{affixes} may be applied to the
\texttt{cs:date-part} elements.

\textbf{Note} Affixes are not allowed on \texttt{cs:date} when defining
localized date formats. This restriction is in place to separate
locale-specific affixes (set on the \texttt{cs:date-part} elements) from
any style-specific affixes (set on the calling \texttt{cs:date}
element), such as parentheses. An example of a macro calling a localized
date format:

\begin{Shaded}
\begin{Highlighting}[]
\NormalTok{\textless{}}\KeywordTok{macro}\OtherTok{ name=}\StringTok{"issued"}\NormalTok{\textgreater{}}
\NormalTok{ \textless{}}\KeywordTok{date}\OtherTok{ variable=}\StringTok{"issued"}\OtherTok{ form=}\StringTok{"numeric"}\OtherTok{ prefix=}\StringTok{"("}\OtherTok{ suffix=}\StringTok{")"}\NormalTok{/\textgreater{}}
\NormalTok{\textless{}/}\KeywordTok{macro}\NormalTok{\textgreater{}}
\end{Highlighting}
\end{Shaded}

\hypertarget{localized-options}{%
\subsubsection{Localized Options}\label{localized-options}}

There are two localized options, \texttt{limit-day-ordinals-to-day-1}
and \texttt{punctuation-in-quote} (see
\protect\hyperlink{locale-options}{Locale Options}). These global
options (which affect both citations and the bibliography) are set as
optional attributes on \texttt{cs:style-options}.

\hypertarget{rendering-elements}{%
\subsection{Rendering Elements}\label{rendering-elements}}

Rendering elements specify which, and in what order, pieces of
bibliographic metadata are included in citations and bibliographies, and
offer control over their formatting.

\hypertarget{layout}{%
\subsubsection{Layout}\label{layout}}

The \texttt{cs:layout} rendering element is a required child element of
\texttt{cs:citation} and \texttt{cs:bibliography}. It must contain one
or more of the other rendering elements described below, and may carry
\protect\hyperlink{affixes}{affixes} and
\protect\hyperlink{formatting}{formatting} attributes. When used within
\texttt{cs:citation}, the \protect\hyperlink{delimiter}{delimiter}
attribute may be used to specify a delimiter for cites within a
citation. For example, a citation like "(1, 2)" can be achieved with:

\begin{Shaded}
\begin{Highlighting}[]
\NormalTok{\textless{}}\KeywordTok{citation}\NormalTok{\textgreater{}}
\NormalTok{  \textless{}}\KeywordTok{layout}\OtherTok{ prefix=}\StringTok{"("}\OtherTok{ suffix=}\StringTok{")"}\OtherTok{ delimiter=}\StringTok{", "}\NormalTok{\textgreater{}}
\NormalTok{    \textless{}}\KeywordTok{text}\OtherTok{ variable=}\StringTok{"citation{-}number"}\NormalTok{/\textgreater{}}
\NormalTok{  \textless{}/}\KeywordTok{layout}\NormalTok{\textgreater{}}
\NormalTok{\textless{}/}\KeywordTok{citation}\NormalTok{\textgreater{}}
\end{Highlighting}
\end{Shaded}

\hypertarget{text}{%
\subsubsection{Text}\label{text}}

The \texttt{cs:text} rendering element outputs text. It must carry one
of the following attributes to select what should be rendered:

\begin{itemize}
\tightlist
\item
  \texttt{variable} - renders the text contents of a variable. Attribute
  value must be one of the
  \protect\hyperlink{standard-variables}{standard variables}. May be
  accompanied by the \texttt{form} attribute to select the "long"
  (default) or "short" form of a variable (e.g. the full or short
  title). If the "short" form is selected but unavailable, the "long"
  form is rendered instead.
\item
  \texttt{macro} - renders the text output of a macro. Attribute value
  must match the value of the \texttt{name} attribute of a
  \texttt{cs:macro} element (see \protect\hyperlink{macro}{Macro}).
\item
  \texttt{term} - renders a term. Attribute value must be one of the
  terms listed in \protect\hyperlink{appendix-ii---terms}{Appendix II -
  Terms}. May be accompanied by the \texttt{plural} attribute to select
  the singular ("false", default) or plural ("true") variant of a term,
  and by the \texttt{form} attribute to select the "long" (default),
  "short", "verb", "verb-short" or "symbol" form variant (see also
  \protect\hyperlink{terms}{Terms}).
\item
  \texttt{value} - renders the attribute value itself.
\end{itemize}

An example of \texttt{cs:text} rendering the "title" variable:

\begin{Shaded}
\begin{Highlighting}[]
\NormalTok{\textless{}}\KeywordTok{text}\OtherTok{ variable=}\StringTok{"title"}\NormalTok{/\textgreater{}}
\end{Highlighting}
\end{Shaded}

\texttt{cs:text} may also carry \protect\hyperlink{affixes}{affixes},
\protect\hyperlink{display}{display},
\protect\hyperlink{formatting}{formatting},
\protect\hyperlink{quotes}{quotes},
\protect\hyperlink{strip-periods}{strip-periods} and
\protect\hyperlink{text-case}{text-case} attributes.

\hypertarget{date}{%
\subsubsection{Date}\label{date}}

The \texttt{cs:date} rendering element outputs the date selected from
the list of \protect\hyperlink{date-variables}{date variables} with the
required \texttt{variable} attribute. A date can be rendered in either a
localized or non-localized format.

\protect\hyperlink{localized-date-formats}{Localized date formats} are
selected with the optional \texttt{form} attribute, which must be set to
either "numeric" (for fully numeric formats, e.g. "12-15-2005"), or
"text" (for formats with a non-numeric month, e.g. "December 15, 2005").
Localized date formats can be customized in two ways. First, the
\texttt{date-parts} attribute may be used to show fewer date parts. The
possible values are:

\begin{itemize}
\tightlist
\item
  "year-month-day" - (default), renders the year, month and day
\item
  "year-month" - renders the year and month
\item
  "year" - renders the year
\end{itemize}

Secondly, \texttt{cs:date} may have one or more \texttt{cs:date-part}
child elements (see \protect\hyperlink{date-part}{Date-part}). The
attributes set on these elements override those specified for the
localized date formats (e.g. to get abbreviated months for all locales,
the \texttt{form} attribute on the month-\texttt{cs:date-part} element
can be set to "short"). These \texttt{cs:date-part} elements do not
affect which, or in what order, date parts are rendered.
\protect\hyperlink{affixes}{Affixes}, which are very locale-specific,
are not allowed on these \texttt{cs:date-part} elements.

In the absence of the \texttt{form} attribute, \texttt{cs:date}
describes a self-contained non-localized date format. In this case, the
date format is constructed using \texttt{cs:date-part} child elements.
With a required \texttt{name} attribute set to either \texttt{day},
\texttt{month} or \texttt{year}, the order of these elements reflects
the display order of respectively the day, month, and year. The date can
be formatted with \protect\hyperlink{formatting}{formatting} attributes
on the \texttt{cs:date-part} elements, as well as several
\texttt{cs:date-part}-specific attributes (see
\protect\hyperlink{date-part}{Date-part}). The
\protect\hyperlink{delimiter}{delimiter} attribute may be set on
\texttt{cs:date} to specify the delimiter for the \texttt{cs:date-part}
elements, and \protect\hyperlink{affixes}{affixes} may be applied to the
\texttt{cs:date-part} elements.

For both localized and non-localized dates, \texttt{cs:date} may carry
\protect\hyperlink{affixes}{affixes},
\protect\hyperlink{display}{display},
\protect\hyperlink{formatting}{formatting} and
\protect\hyperlink{text-case}{text-case} attributes.

\hypertarget{date-part}{%
\paragraph{Date-part}\label{date-part}}

The \texttt{cs:date-part} elements control how date parts are rendered.
Unless the parent \texttt{cs:date} element calls a localized date
format, they also determine which, and in what order, date parts appear.
A \texttt{cs:date-part} element describes the date part selected with
the required \texttt{name} attribute:

\begin{description}
\item["day"]
For "day", \texttt{cs:date-part} may carry the \texttt{form} attribute,
with values:

\begin{itemize}
\tightlist
\item
  "numeric" - (default), e.g. "1"
\item
  "numeric-leading-zeros" - e.g. "01"
\item
  "ordinal" - e.g. "1st"
\end{itemize}

Some languages, such as French, only use the "ordinal" form for the
first day of the month ("1er janvier", "2 janvier", "3 janvier", etc.).
Such output can be achieved with the "ordinal" form and use of the
\texttt{limit-day-ordinals-to-day-1} attribute (see
\protect\hyperlink{locale-options}{Locale Options}).
\item["month"]
For "month", \texttt{cs:date-part} may carry the
\protect\hyperlink{strip-periods}{strip-periods} and \texttt{form}
attributes. In locale files, month abbreviations (the "short" form of
the month \protect\hyperlink{terms}{terms}) should be defined with
periods if applicable (e.g. "Jan.", "Feb.", etc.). These periods can be
removed by setting \protect\hyperlink{strip-periods}{strip-periods} to
"true" ("false" is the default). The \texttt{form} attribute can be set
to:

\begin{itemize}
\tightlist
\item
  "long" - (default), e.g. "January"
\item
  "short" - e.g. "Jan."
\item
  "numeric" - e.g. "1"
\item
  "numeric-leading-zeros" - e.g. "01"
\end{itemize}
\item["year"]
For "year", \texttt{cs:date-part} may carry the \texttt{form} attribute,
with values:

\begin{itemize}
\tightlist
\item
  "long" - (default), e.g. "2005"
\item
  "short" - e.g. "05"
\end{itemize}
\end{description}

\texttt{cs:date-part} may also carry
\protect\hyperlink{formatting}{formatting},
\protect\hyperlink{text-case}{text-case} and \texttt{range-delimiter}
(see \protect\hyperlink{date-ranges}{Date Ranges}) attributes.
Attributes for \protect\hyperlink{affixes}{affixes} are allowed, unless
\texttt{cs:date} calls a localized date format.

\hypertarget{date-ranges}{%
\subparagraph{Date Ranges}\label{date-ranges}}

The default delimiter for dates in a date range is an en-dash (e.g. "May
-- July 2008"). Custom range delimiters can be set on
\texttt{cs:date-part} elements with the optional
\texttt{range-delimiter} attribute. When a date range is rendered, the
range delimiter is drawn from the \texttt{cs:date-part} element matching
the largest date part ("year", "month", or "day") that differs between
the two dates. For example,

\begin{Shaded}
\begin{Highlighting}[]
\NormalTok{\textless{}}\KeywordTok{style}\NormalTok{\textgreater{}}
\NormalTok{  \textless{}}\KeywordTok{citation}\NormalTok{\textgreater{}}
\NormalTok{    \textless{}}\KeywordTok{layout}\NormalTok{\textgreater{}}
\NormalTok{      \textless{}}\KeywordTok{date}\OtherTok{ variable=}\StringTok{"issued"}\NormalTok{\textgreater{}}
\NormalTok{        \textless{}}\KeywordTok{date{-}part}\OtherTok{ name=}\StringTok{"day"}\OtherTok{ suffix=}\StringTok{" "}\OtherTok{ range{-}delimiter=}\StringTok{"{-}"}\NormalTok{/\textgreater{}}
\NormalTok{        \textless{}}\KeywordTok{date{-}part}\OtherTok{ name=}\StringTok{"month"}\OtherTok{ suffix=}\StringTok{" "}\NormalTok{/\textgreater{}}
\NormalTok{        \textless{}}\KeywordTok{date{-}part}\OtherTok{ name=}\StringTok{"year"}\OtherTok{ range{-}delimiter=}\StringTok{"/"}\NormalTok{/\textgreater{}}
\NormalTok{      \textless{}/}\KeywordTok{date}\NormalTok{\textgreater{}}
\NormalTok{    \textless{}/}\KeywordTok{layout}\NormalTok{\textgreater{}}
\NormalTok{  \textless{}/}\KeywordTok{citation}\NormalTok{\textgreater{}}
\NormalTok{\textless{}/}\KeywordTok{style}\NormalTok{\textgreater{}}
\end{Highlighting}
\end{Shaded}

would result in "1-4 May 2008", "May -- July 2008" and "May 2008/June
2009".

\hypertarget{ad-and-bc}{%
\paragraph{AD and BC}\label{ad-and-bc}}

The "ad" term (Anno Domini) is automatically appended to positive years
of less than four digits (e.g. "79" becomes "79AD"). The "bc" term
(Before Christ) is automatically appended to negative years (e.g.
"-2500" becomes "2500BC").

\hypertarget{seasons}{%
\paragraph{Seasons}\label{seasons}}

If a date includes a season instead of a month, a season term
("season-01" to "season-04", respectively Spring, Summer, Autumn and
Winter) take the place of the month term. E.g.,

\begin{Shaded}
\begin{Highlighting}[]
\NormalTok{\textless{}}\KeywordTok{style}\NormalTok{\textgreater{}}
\NormalTok{  \textless{}}\KeywordTok{citation}\NormalTok{\textgreater{}}
\NormalTok{    \textless{}}\KeywordTok{layout}\NormalTok{\textgreater{}}
\NormalTok{      \textless{}}\KeywordTok{date}\OtherTok{ variable=}\StringTok{"issued"}\NormalTok{\textgreater{}}
\NormalTok{        \textless{}}\KeywordTok{date{-}part}\OtherTok{ name=}\StringTok{"month"}\OtherTok{ suffix=}\StringTok{" "}\NormalTok{/\textgreater{}}
\NormalTok{        \textless{}}\KeywordTok{date{-}part}\OtherTok{ name=}\StringTok{"year"}\NormalTok{/\textgreater{}}
\NormalTok{      \textless{}/}\KeywordTok{date}\NormalTok{\textgreater{}}
\NormalTok{    \textless{}/}\KeywordTok{layout}\NormalTok{\textgreater{}}
\NormalTok{  \textless{}/}\KeywordTok{citation}\NormalTok{\textgreater{}}
\NormalTok{\textless{}/}\KeywordTok{style}\NormalTok{\textgreater{}}
\end{Highlighting}
\end{Shaded}

would result in "May 2008" and "Winter 2009".

\hypertarget{approximate-dates}{%
\paragraph{Approximate Dates}\label{approximate-dates}}

Approximate dates test "true" for the \texttt{is-uncertain-date}
conditional (see \protect\hyperlink{choose}{Choose}). For example,

\begin{Shaded}
\begin{Highlighting}[]
\NormalTok{\textless{}}\KeywordTok{style}\NormalTok{\textgreater{}}
\NormalTok{  \textless{}}\KeywordTok{citation}\NormalTok{\textgreater{}}
\NormalTok{    \textless{}}\KeywordTok{layout}\NormalTok{\textgreater{}}
\NormalTok{      \textless{}}\KeywordTok{choose}\NormalTok{\textgreater{}}
\NormalTok{        \textless{}}\KeywordTok{if}\OtherTok{ is{-}uncertain{-}date=}\StringTok{"issued"}\NormalTok{\textgreater{}}
\NormalTok{          \textless{}}\KeywordTok{text}\OtherTok{ term=}\StringTok{"circa"}\OtherTok{ form=}\StringTok{"short"}\OtherTok{ suffix=}\StringTok{" "}\NormalTok{/\textgreater{}}
\NormalTok{        \textless{}/}\KeywordTok{if}\NormalTok{\textgreater{}}
\NormalTok{      \textless{}/}\KeywordTok{choose}\NormalTok{\textgreater{}}
\NormalTok{      \textless{}}\KeywordTok{date}\OtherTok{ variable=}\StringTok{"issued"}\NormalTok{\textgreater{}}
\NormalTok{        \textless{}}\KeywordTok{date{-}part}\OtherTok{ name=}\StringTok{"year"}\NormalTok{/\textgreater{}}
\NormalTok{      \textless{}/}\KeywordTok{date}\NormalTok{\textgreater{}}
\NormalTok{    \textless{}/}\KeywordTok{layout}\NormalTok{\textgreater{}}
\NormalTok{  \textless{}/}\KeywordTok{citation}\NormalTok{\textgreater{}}
\NormalTok{\textless{}/}\KeywordTok{style}\NormalTok{\textgreater{}}
\end{Highlighting}
\end{Shaded}

would result in "2005" (normal date) and "ca. 2003" (approximate date).

\hypertarget{number}{%
\subsubsection{Number}\label{number}}

The \texttt{cs:number} rendering element outputs the number variable
selected with the required \texttt{variable} attribute.
\protect\hyperlink{number-variables}{Number variables} are a subset of
the list of \protect\hyperlink{standard-variables}{standard variables}.

If a number variable is rendered with \texttt{cs:number} and only
contains numeric content (as determined by the rules for
\texttt{is-numeric}, see \protect\hyperlink{choose}{Choose}), the
number(s) are extracted. Variable content is rendered "as is" when the
variable contains any non-numeric content (e.g. "Special edition").

During the extraction, numbers separated by a hyphen are stripped of
intervening spaces ("2 - 4" becomes "2-4"). Numbers separated by a comma
receive one space after the comma ("2,3" and "2 , 3" become "2, 3"),
while numbers separated by an ampersand receive one space before and one
after the ampersand ("2\&3" becomes "2 \& 3").

Extracted numbers can be formatted via the optional \texttt{form}
attribute, with values:

\begin{itemize}
\tightlist
\item
  "numeric" - (default), e.g. "1", "2", "3"
\item
  "ordinal" - e.g. "1st", "2nd", "3rd". Ordinal suffixes are defined
  with terms (see \protect\hyperlink{ordinal-suffixes}{Ordinal
  Suffixes}).
\item
  "long-ordinal" - e.g. "first", "second", "third". Long ordinals are
  defined with the \protect\hyperlink{terms}{terms} "long-ordinal-01" to
  "long-ordinal-10", which are used for the numbers 1 through 10. For
  other numbers "long-ordinal" falls back to "ordinal".
\item
  "roman" - e.g. "i", "ii", "iii"
\end{itemize}

Numbers with prefixes or suffixes are never ordinalized or rendered in
roman numerals (e.g. "2E" remains "2E). Numbers without affixes are
individually transformed ("2, 3" can become "2nd, 3rd", "second, third"
or "ii, iii").

\texttt{cs:number} may carry \protect\hyperlink{affixes}{affixes},
\protect\hyperlink{display}{display},
\protect\hyperlink{formatting}{formatting} and
\protect\hyperlink{text-case}{text-case} attributes.

\hypertarget{names}{%
\subsubsection{Names}\label{names}}

The \texttt{cs:names} rendering element outputs the contents of one or
more \protect\hyperlink{name-variables}{name variables} (selected with
the required \texttt{variable} attribute), each of which can contain
multiple names (e.g. the "author" variable contains all the author names
of the cited item). If multiple variables are selected (separated by
single spaces, see example below), each variable is independently
rendered in the order specified, with one exception: when the selection
consists of "editor" and "translator", and when the contents of these
two name variables is identical, then the contents of only one name
variable is rendered. In addition, the "editortranslator" term is used
if the \texttt{cs:names} element contains a \texttt{cs:label} element,
replacing the default "editor" and "translator" terms (e.g. resulting in
"Doe (editor \& translator)"). The
\protect\hyperlink{delimiter}{delimiter} attribute may be set on
\texttt{cs:names} to separate the names of the different name variables
(e.g. the semicolon in "Doe, Smith (editors); Johnson (translator)").

\begin{Shaded}
\begin{Highlighting}[]
\NormalTok{\textless{}}\KeywordTok{names}\OtherTok{ variable=}\StringTok{"editor translator"}\OtherTok{ delimiter=}\StringTok{"; "}\NormalTok{\textgreater{}}
\NormalTok{  \textless{}}\KeywordTok{label}\OtherTok{ prefix=}\StringTok{" ("}\OtherTok{ suffix=}\StringTok{")"}\NormalTok{/\textgreater{}}
\NormalTok{\textless{}/}\KeywordTok{names}\NormalTok{\textgreater{}}
\end{Highlighting}
\end{Shaded}

\texttt{cs:names} has four child elements (discussed below):
\texttt{cs:name}, \texttt{cs:et-al}, \texttt{cs:substitute} and
\texttt{cs:label}. The \texttt{cs:names} element may carry
\protect\hyperlink{affixes}{affixes},
\protect\hyperlink{display}{display} and
\protect\hyperlink{formatting}{formatting} attributes.

\hypertarget{name}{%
\paragraph{Name}\label{name}}

The \texttt{cs:name} element, an optional child element of
\texttt{cs:names}, can be used to describe the formatting of individual
names, and the separation of names within a name variable.
\texttt{cs:name} may carry the following attributes:

\begin{description}
\item[\texttt{and}]
Specifies the delimiter between the second to last and last name of the
names in a name variable. Allowed values are "text" (selects the "and"
term, e.g. "Doe, Johnson and Smith") and "symbol" (selects the
ampersand, e.g. "Doe, Johnson \& Smith").
\item[\texttt{delimiter}]
Specifies the text string used to separate names in a name variable.
Default is ", " (e.g. "Doe, Smith").
\item[\texttt{delimiter-precedes-et-al}]
Determines when the name delimiter or a space is used between a
truncated name list and the "et-al" (or "and others") term in case of
et-al abbreviation. Allowed values:

\begin{itemize}
\tightlist
\item
  "contextual" - (default), name delimiter is only used for name lists
  truncated to two or more names

  \begin{itemize}
  \tightlist
  \item
    1 name: "J. Doe et al."
  \item
    2 names: "J. Doe, S. Smith, et al."
  \end{itemize}
\item
  "after-inverted-name" - name delimiter is only used if the preceding
  name is inverted as a result of the \texttt{name-as-sort-order}
  attribute. E.g. with \texttt{name-as-sort-order} set to "first":

  \begin{itemize}
  \tightlist
  \item
    "Doe, J., et al."
  \item
    "Doe, J., S. Smith et al."
  \end{itemize}
\item
  "always" - name delimiter is always used

  \begin{itemize}
  \tightlist
  \item
    1 name: "J. Doe, et al."
  \item
    2 names: "J. Doe, S. Smith, et al."
  \end{itemize}
\item
  "never" - name delimiter is never used

  \begin{itemize}
  \tightlist
  \item
    1 name: "J. Doe et al."
  \item
    2 names: "J. Doe, S. Smith et al."
  \end{itemize}
\end{itemize}
\item[\texttt{delimiter-precedes-last}]
Determines when the name delimiter is used to separate the second to
last and the last name in name lists (if \texttt{and} is not set, the
name delimiter is always used, regardless of the value of
\texttt{delimiter-precedes-last}). Allowed values:

\begin{itemize}
\tightlist
\item
  "contextual" - (default), name delimiter is only used for name lists
  with three or more names

  \begin{itemize}
  \tightlist
  \item
    2 names: "J. Doe and T. Williams"
  \item
    3 names: "J. Doe, S. Smith, and T. Williams"
  \end{itemize}
\item
  "after-inverted-name" - name delimiter is only used if the preceding
  name is inverted as a result of the \texttt{name-as-sort-order}
  attribute. E.g. with \texttt{name-as-sort-order} set to "first":

  \begin{itemize}
  \tightlist
  \item
    "Doe, J., and T. Williams"
  \item
    "Doe, J., S. Smith and T. Williams"
  \end{itemize}
\item
  "always" - name delimiter is always used

  \begin{itemize}
  \tightlist
  \item
    2 names: "J. Doe, and T. Williams"
  \item
    3 names: "J. Doe, S. Smith, and T. Williams"
  \end{itemize}
\item
  "never" - name delimiter is never used

  \begin{itemize}
  \tightlist
  \item
    2 names: "J. Doe and T. Williams"
  \item
    3 names: "J. Doe, S. Smith and T. Williams"
  \end{itemize}
\end{itemize}
\item[\texttt{et-al-min} / \texttt{et-al-use-first}]
Use of these two attributes enables et-al abbreviation. If the number of
names in a name variable matches or exceeds the number set on
\texttt{et-al-min}, the rendered name list is truncated after reaching
the number of names set on \texttt{et-al-use-first}. The "et-al" (or
"and others") term is appended to truncated name lists (see also
\protect\hyperlink{et-al}{Et-al}). By default, when a name list is
truncated to a single name, the name and the "et-al" (or "and others")
term are separated by a space (e.g. "Doe et al."). When a name list is
truncated to two or more names, the name delimiter is used (e.g. "Doe,
Smith, et al."). This behavior can be changed with the
\texttt{delimiter-precedes-et-al} attribute.
\item[\texttt{et-al-subsequent-min} /
\texttt{et-al-subsequent-use-first}]
If used, the values of these attributes replace those of respectively
\texttt{et-al-min} and \texttt{et-al-use-first} for subsequent cites
(cites referencing earlier cited items).
\item[\texttt{et-al-use-last}]
When set to "true" (the default is "false"), name lists truncated by
et-al abbreviation are followed by the name delimiter, the ellipsis
character, and the last name of the original name list. This is only
possible when the original name list has at least two more names than
the truncated name list (for this the value of
\texttt{et-al-use-first}/\texttt{et-al-subsequent-min} must be at least
2 less than the value of
\texttt{et-al-min}/\texttt{et-al-subsequent-use-first}). An example:

\begin{verbatim}
A. Goffeau, B. G. Barrell, H. Bussey, R. W. Davis, B. Dujon, H.
Feldmann, … S. G. Oliver
\end{verbatim}
\end{description}

The remaining attributes, discussed below, only affect personal names.
Personal names require a "family" name-part, and may also contain
"given", "suffix", "non-dropping-particle" and "dropping-particle"
name-parts. These name-parts are defined as:

\begin{itemize}
\tightlist
\item
  "family" - surname minus any particles and suffixes
\item
  "given" - given names, either full ("John Edward") or initialized ("J.
  E.")
\item
  "suffix" - name suffix, e.g. "Jr." in "John Smith Jr." and "III" in
  "Bill Gates III"
\item
  "non-dropping-particle" - name particles that are not dropped when
  only the surname is shown ("van" in the Dutch surname "van Gogh") but
  which may be treated separately from the family name, e.g. for sorting
\item
  "dropping-particle" - name particles that are dropped when only the
  surname is shown ("van" in "Ludwig van Beethoven", which becomes
  "Beethoven", or "von" in "Alexander von Humboldt", which becomes
  "Humboldt")
\end{itemize}

The attributes affecting personal names:

\begin{description}
\item[\texttt{form}]
Specifies whether all the name-parts of personal names should be
displayed (value "long", the default), or only the family name and the
non-dropping-particle (value "short"). A third value, "count", returns
the total number of names that would otherwise be rendered by the use of
the \texttt{cs:names} element (taking into account the effects of et-al
abbreviation and editor/translator collapsing), which allows for
advanced \protect\hyperlink{sorting}{sorting}.
\item[\texttt{initialize}]
When set to "false" (the default is "true"), given names are no longer
initialized when "initialize-with" is set. However, the value of
"initialize-with" is still added after initials present in the full name
(e.g. with \texttt{initialize} set to "false", and
\texttt{initialize-with} set to ".", "James T Kirk" becomes "James T.
Kirk").
\item[\texttt{initialize-with}]
When set, given names are converted to initials. The attribute value is
added after each initial ("." results in "J. J. Doe"). For compound
given names (e.g. "Jean-Luc"), hyphenation of the initials can be
controlled with the global \texttt{initialize-with-hyphen} option (see
\protect\hyperlink{hyphenation-of-initialized-names}{Hyphenation of
Initialized Names}).
\item[\texttt{name-as-sort-order}]
Specifies that names should be displayed with the given name following
the family name (e.g. "John Doe" becomes "Doe, John"). The attribute has
two possible values:

\begin{itemize}
\tightlist
\item
  "first" - attribute only has an effect on the first name of each name
  variable
\item
  "all" - attribute has an effect on all names
\end{itemize}

Note that even when \texttt{name-as-sort-order} changes the name-part
order, the display order is not necessarily the same as the sorting
order for names containing particles and suffixes (see
\protect\hyperlink{name-part-order}{Name-part order}). Also,
\texttt{name-as-sort-order} only affects names written in scripts where
the given name typically precedes the family name, such as Latin, Greek,
Cyrillic and Arabic. In contrast, names written in Asian scripts are
always displayed with the family name preceding the given name.
\item[\texttt{sort-separator}]
Sets the delimiter for name-parts that have switched positions as a
result of \texttt{name-as-sort-order}. The default value is ", " ("Doe,
John"). As is the case for \texttt{name-as-sort-order}, this attribute
only affects names in scripts that know "given-name family-name" order.
\end{description}

\texttt{cs:name} may also carry \protect\hyperlink{affixes}{affixes} and
\protect\hyperlink{formatting}{formatting} attributes.

\hypertarget{name-part-order}{%
\subparagraph{Name-part Order}\label{name-part-order}}

The order of name-parts depends on the values of the \texttt{form} and
\texttt{name-as-sort-order} attributes on \texttt{cs:name}, the value of
the \texttt{demote-non-dropping-particle} attribute on \texttt{cs:style}
(one of the \protect\hyperlink{global-options}{global options}), and the
script of the individual name. Note that the display and sorting order
of name-parts often differs. An overview of the possible orders:

\textbf{Display order of names in "given-name family-name" scripts
(Latin, etc.)}

\begin{center}\rule{0.5\linewidth}{0.5pt}\end{center}

\begin{description}
\item[Conditions]
\texttt{form} set to "long"
\item[Order]
\begin{enumerate}
\def\labelenumi{\arabic{enumi})}
\tightlist
\item[]
\item
  given
\item
  dropping-particle
\item
  non-dropping-particle
\item
  family
\item
  suffix
\end{enumerate}
\item[Example]
{[}Vincent{]} {[}{]} {[}van{]} {[}Gogh{]} {[}III{]}
\item[Example]
{[}Alexander{]} {[}von{]} {[}{]} {[}Humboldt{]} {[}Jr.{]}
\end{description}

\begin{center}\rule{0.5\linewidth}{0.5pt}\end{center}

\begin{description}
\item[Conditions]
\texttt{form} set to "long", name-as-sort-order active,
\texttt{demote-non-dropping-particle} set to "never" or "sort-only"
\item[Order]
\begin{enumerate}
\def\labelenumi{\arabic{enumi})}
\tightlist
\item[]
\item
  non-dropping-particle
\item
  family
\item
  given
\item
  dropping-particle
\item
  suffix
\end{enumerate}
\item[Example]
{[}van{]} {[}Gogh{]}, {[}Vincent{]} {[}{]}, {[}III{]}
\end{description}

\begin{center}\rule{0.5\linewidth}{0.5pt}\end{center}

\begin{description}
\item[Conditions]
\texttt{form} set to "long", name-as-sort-order active,
\texttt{demote-non-dropping-particle} set to "display-and-sort"
\item[Order]
\begin{enumerate}
\def\labelenumi{\arabic{enumi})}
\tightlist
\item[]
\item
  family
\item
  given
\item
  dropping-particle
\item
  non-dropping-particle
\item
  suffix
\end{enumerate}
\item[Example]
{[}Gogh{]}, {[}Vincent{]} {[}{]} {[}van{]}, {[}III{]}
\end{description}

\begin{center}\rule{0.5\linewidth}{0.5pt}\end{center}

\begin{description}
\item[Conditions]
\texttt{form} set to "short"
\item[Order]
\begin{enumerate}
\def\labelenumi{\arabic{enumi})}
\tightlist
\item[]
\item
  non-dropping-particles
\item
  family
\end{enumerate}
\item[Example]
{[}van{]} {[}Gogh{]}
\end{description}

\begin{center}\rule{0.5\linewidth}{0.5pt}\end{center}

\textbf{Sorting order of names in "given-name family-name" scripts
(Latin, etc.)}

N.B. The sort keys are listed in descending order of priority.

\begin{center}\rule{0.5\linewidth}{0.5pt}\end{center}

\begin{description}
\item[Conditions]
\texttt{demote-non-dropping-particle} set to "never"
\item[Order]
\begin{enumerate}
\def\labelenumi{\arabic{enumi})}
\tightlist
\item[]
\item
  non-dropping-particle + family
\item
  dropping-particle
\item
  given
\item
  suffix
\end{enumerate}
\item[Example]
{[}van Gogh{]} {[}{]} {[}Vincent{]} {[}III{]}
\end{description}

\begin{center}\rule{0.5\linewidth}{0.5pt}\end{center}

\begin{description}
\item[Conditions]
\texttt{demote-non-dropping-particle} set to "sort-only" or
"display-and-sort"
\item[Order]
\begin{enumerate}
\def\labelenumi{\arabic{enumi})}
\tightlist
\item[]
\item
  family
\item
  dropping-particle + non-dropping-particle
\item
  given
\item
  suffix
\end{enumerate}
\item[Example]
{[}Gogh{]} {[}van{]} {[}Vincent{]} {[}III{]}
\end{description}

\begin{center}\rule{0.5\linewidth}{0.5pt}\end{center}

\textbf{Display and sorting order of names in "family-name given-name"
scripts (Chinese, etc.)}

\begin{center}\rule{0.5\linewidth}{0.5pt}\end{center}

\begin{description}
\item[Conditions]
\texttt{form} set to "long"
\item[Order]
\begin{enumerate}
\def\labelenumi{\arabic{enumi})}
\tightlist
\item[]
\item
  family
\item
  given
\end{enumerate}
\item[Example]
毛 泽 东 {[}Mao Zedong{]}
\end{description}

\begin{center}\rule{0.5\linewidth}{0.5pt}\end{center}

\begin{description}
\item[Conditions]
\texttt{form} set to "short"
\item[Order]
\begin{enumerate}
\def\labelenumi{\arabic{enumi})}
\tightlist
\item[]
\item
  family
\end{enumerate}
\item[Example]
毛 {[}Mao{]}
\end{description}

\begin{center}\rule{0.5\linewidth}{0.5pt}\end{center}

Non-personal names lack name-parts and are sorted as is, although
English articles ("a", "an" and "the") at the start of the name are
stripped. For example, "The New York Times" sorts as "New York Times".

\hypertarget{name-part-formatting}{%
\subparagraph{Name-part Formatting}\label{name-part-formatting}}

The \texttt{cs:name} element may contain one or two
\texttt{cs:name-part} child elements for name-part-specific formatting.
\texttt{cs:name-part} must carry the \texttt{name} attribute, set to
either "given" or "family".

If set to "given", \protect\hyperlink{formatting}{formatting} and
\protect\hyperlink{text-case}{text-case} attributes on
\texttt{cs:name-part} affect the "given" and "dropping-particle"
name-parts. \protect\hyperlink{affixes}{affixes} surround the "given"
name-part, enclosing any demoted name particles for inverted names.

If set to "family", \protect\hyperlink{formatting}{formatting} and
\protect\hyperlink{text-case}{text-case} attributes affect the "family"
and "non-dropping-particle" name-parts.
\protect\hyperlink{affixes}{affixes} surround the "family" name-part,
enclosing any preceding name particles, as well as the "suffix"
name-part for non-inverted names.

The "suffix" name-part is not subject to name-part formatting. The use
of \texttt{cs:name-part} elements does not influence which, or in what
order, name-parts are rendered. An example, yielding names like "Jane
DOE":

\begin{Shaded}
\begin{Highlighting}[]
\NormalTok{\textless{}}\KeywordTok{names}\OtherTok{ variable=}\StringTok{"author"}\NormalTok{\textgreater{}}
\NormalTok{  \textless{}}\KeywordTok{name}\NormalTok{\textgreater{}}
\NormalTok{    \textless{}}\KeywordTok{name{-}part}\OtherTok{ name=}\StringTok{"family"}\OtherTok{ text{-}case=}\StringTok{"uppercase"}\NormalTok{/\textgreater{}}
\NormalTok{  \textless{}/}\KeywordTok{name}\NormalTok{\textgreater{}}
\NormalTok{\textless{}/}\KeywordTok{names}\NormalTok{\textgreater{}}
\end{Highlighting}
\end{Shaded}

\hypertarget{et-al}{%
\paragraph{Et-al}\label{et-al}}

Et-al abbreviation, controlled via the \texttt{et-al-…} attributes (see
\protect\hyperlink{name}{Name}), can be further customized with the
optional \texttt{cs:et-al} element, which must follow the
\texttt{cs:name} element (if present).

The \protect\hyperlink{formatting}{formatting} attributes may be used on
\texttt{cs:et-al}, for example to italicize the "et-al" term:

\begin{Shaded}
\begin{Highlighting}[]
\NormalTok{\textless{}}\KeywordTok{names}\OtherTok{ variable=}\StringTok{"author"}\NormalTok{\textgreater{}}
\NormalTok{  \textless{}}\KeywordTok{et{-}al}\OtherTok{ font{-}style=}\StringTok{"italic"}\NormalTok{/\textgreater{}}
\NormalTok{\textless{}/}\KeywordTok{names}\NormalTok{\textgreater{}}
\end{Highlighting}
\end{Shaded}

The \texttt{term} attribute may also be set, to either "et-al" (the
default) or "and others", to use either term:

\begin{Shaded}
\begin{Highlighting}[]
\NormalTok{\textless{}}\KeywordTok{names}\OtherTok{ variable=}\StringTok{"author"}\NormalTok{\textgreater{}}
\NormalTok{  \textless{}}\KeywordTok{et{-}al}\OtherTok{ term=}\StringTok{"and others"}\NormalTok{/\textgreater{}}
\NormalTok{\textless{}/}\KeywordTok{names}\NormalTok{\textgreater{}}
\end{Highlighting}
\end{Shaded}

\hypertarget{substitute}{%
\paragraph{Substitute}\label{substitute}}

The optional \texttt{cs:substitute} element, which must be included as
the last child element of \texttt{cs:names}, adds substitution in case
the \protect\hyperlink{name-variables}{name variables} specified in the
parent \texttt{cs:names} element are empty. The substitutions are
specified as child elements of \texttt{cs:substitute}, and must consist
of one or more \protect\hyperlink{rendering-elements}{rendering
elements} (with the exception of \texttt{cs:layout}). A shorthand
version of \texttt{cs:names} without child elements, which inherits the
attributes values set on the \texttt{cs:name} and \texttt{cs:et-al}
child elements of the original \texttt{cs:names} element, may also be
used. If \texttt{cs:substitute} contains multiple child elements, the
first element to return a non-empty result is used for substitution.
Substituted variables are suppressed in the rest of the output to
prevent duplication. Substituted variables are considered empty for the
purposes of determining whether to suppress an enclosing
\texttt{cs:group}. If the variable was rendered earlier in the citation,
before the "substitute" element, it is not suppressed. An example, where
an empty "author" name variable is substituted by the "editor" name
variable, or, when no editors exist, by the "title" macro:

\begin{Shaded}
\begin{Highlighting}[]
\NormalTok{\textless{}}\KeywordTok{macro}\OtherTok{ name=}\StringTok{"author"}\NormalTok{\textgreater{}}
\NormalTok{  \textless{}}\KeywordTok{names}\OtherTok{ variable=}\StringTok{"author"}\NormalTok{\textgreater{}}
\NormalTok{    \textless{}}\KeywordTok{substitute}\NormalTok{\textgreater{}}
\NormalTok{      \textless{}}\KeywordTok{names}\OtherTok{ variable=}\StringTok{"editor"}\NormalTok{/\textgreater{}}
\NormalTok{      \textless{}}\KeywordTok{text}\OtherTok{ macro=}\StringTok{"title"}\NormalTok{/\textgreater{}}
\NormalTok{    \textless{}/}\KeywordTok{substitute}\NormalTok{\textgreater{}}
\NormalTok{  \textless{}/}\KeywordTok{names}\NormalTok{\textgreater{}}
\NormalTok{\textless{}/}\KeywordTok{macro}\NormalTok{\textgreater{}}
\end{Highlighting}
\end{Shaded}

\hypertarget{label-in-csnames}{%
\paragraph{\texorpdfstring{Label in
\texttt{cs:names}}{Label in cs:names}}\label{label-in-csnames}}

A \texttt{cs:label} element (see \protect\hyperlink{label}{label}) may
optionally be included in \texttt{cs:names}. It must appear before the
\texttt{cs:substitute} element. The position of \texttt{cs:label}
relative to \texttt{cs:name} determines the order of the name and label
in the rendered text. When used as a child element of \texttt{cs:names},
\texttt{cs:label} does not carry the \texttt{variable} attribute; it
uses the variable(s) set on the parent \texttt{cs:names} element
instead. A second difference is that the \texttt{form} attribute may
also be set to "verb" or "verb-short", so that the allowed values are:

\begin{itemize}
\tightlist
\item
  "long" - (default), e.g. "editor" and "editors" for the "editor" term
\item
  "short" - e.g. "ed." and "eds." for the term "editor"
\item
  "verb" - e.g. "edited by" for the term "editor"
\item
  "verb-short" - e.g. "ed." for the term "editor"
\item
  "symbol" - e.g. "§" and "§§" for the term "section"
\end{itemize}

\hypertarget{label}{%
\subsubsection{Label}\label{label}}

The \texttt{cs:label} rendering element outputs the term matching the
variable selected with the required \texttt{variable} attribute, which
must be set to "locator", "page", or one of the
\protect\hyperlink{number-variables}{number variables}. The term is only
rendered if the selected variable is non-empty. For example,

\begin{Shaded}
\begin{Highlighting}[]
\NormalTok{\textless{}}\KeywordTok{group}\OtherTok{ delimiter=}\StringTok{" "}\NormalTok{\textgreater{}}
\NormalTok{  \textless{}}\KeywordTok{label}\OtherTok{ variable=}\StringTok{"page"}\NormalTok{/\textgreater{}}
\NormalTok{  \textless{}}\KeywordTok{text}\OtherTok{ variable=}\StringTok{"page"}\NormalTok{/\textgreater{}}
\NormalTok{\textless{}/}\KeywordTok{group}\NormalTok{\textgreater{}}
\end{Highlighting}
\end{Shaded}

can result in "page 3" or "pages 5-7". \texttt{cs:label} may carry the
following attributes:

\begin{description}
\item[\texttt{form}]
Selects the form of the term, with allowed values:

\begin{itemize}
\tightlist
\item
  "long" - (default), e.g. "page"/"pages" for the "page" term
\item
  "short" - e.g. "p."/"pp." for the "page" term
\item
  "symbol" - e.g. "§"/"§§" for the "section" term
\end{itemize}
\item[\texttt{plural}]
Sets pluralization of the term, with allowed values:

\begin{itemize}
\tightlist
\item
  "contextual" - (default), the term plurality matches that of the
  variable content. Content is considered plural when it contains
  multiple numbers (e.g. "page 1", "pages 1-3", "volume 2", "volumes 2
  \& 4"), or, in the case of the "number-of-pages" and
  "number-of-volumes" variables, when the number is higher than 1 ("1
  volume" and "3 volumes").
\item
  "always" - always use the plural form, e.g. "pages 1" and "pages 1-3"
\item
  "never" - always use the singular form, e.g. "page 1" and "page 1-3"
\end{itemize}
\end{description}

\texttt{cs:label} may also carry \protect\hyperlink{affixes}{affixes},
\protect\hyperlink{formatting}{formatting},
\protect\hyperlink{text-case}{text-case} and
\protect\hyperlink{strip-periods}{strip-periods} attributes.

\hypertarget{group}{%
\subsubsection{Group}\label{group}}

The \texttt{cs:group} rendering element must contain one or more
\protect\hyperlink{rendering-elements}{rendering elements} (with the
exception of \texttt{cs:layout}). \texttt{cs:group} may carry the
\protect\hyperlink{delimiter}{delimiter} attribute to separate its child
elements, as well as \protect\hyperlink{affixes}{affixes},
\protect\hyperlink{display}{display}, and
\protect\hyperlink{formatting}{formatting} attributes (applied to the
output of the group as a whole). \texttt{cs:group} implicitly acts as a
conditional: \texttt{cs:group} and its child elements are suppressed if
a) at least one rendering element in \texttt{cs:group} calls a variable
(either directly or via a macro), and b) all variables that are called
are empty. This accommodates descriptive {cs:text} and
`cs:label`elements. For example,

\begin{Shaded}
\begin{Highlighting}[]
\NormalTok{\textless{}}\KeywordTok{layout}\NormalTok{\textgreater{}}
\NormalTok{  \textless{}}\KeywordTok{group}\OtherTok{ delimiter=}\StringTok{" "}\NormalTok{\textgreater{}}
\NormalTok{    \textless{}}\KeywordTok{text}\OtherTok{ term=}\StringTok{"retrieved"}\NormalTok{/\textgreater{}}
\NormalTok{    \textless{}}\KeywordTok{text}\OtherTok{ term=}\StringTok{"from"}\NormalTok{/\textgreater{}}
\NormalTok{    \textless{}}\KeywordTok{text}\OtherTok{ variable=}\StringTok{"URL"}\NormalTok{/\textgreater{}}
\NormalTok{  \textless{}/}\KeywordTok{group}\NormalTok{\textgreater{}}
\NormalTok{\textless{}/}\KeywordTok{layout}\NormalTok{\textgreater{}}
\end{Highlighting}
\end{Shaded}

can result in "retrieved from
\url{https://doi.org/10.1128/AEM.02591-07}", but doesn\textquotesingle t
generate output when the "URL" variable is empty.

If a \texttt{cs:group} is nested within another \texttt{cs:group}, the
inner group is evaluated first: a non-empty nested \texttt{cs:group} is
treated as a non-empty variable for the puropses of determining
suppression of the outer \texttt{cs:group}.

When a \texttt{cs:group} contains a child \texttt{cs:macro}, if the
\texttt{cs:macro} is non-empty, it is treated as a non-empty variable
for the purposes of determining suppression of the outer
\texttt{cs:group}.

\hypertarget{choose}{%
\subsubsection{Choose}\label{choose}}

The \texttt{cs:choose} rendering element allows for conditional
rendering of \protect\hyperlink{rendering-elements}{rendering elements}.
An example that renders the "issued" date variable when it exists, and
the "no date" term when it doesn\textquotesingle t:

\begin{Shaded}
\begin{Highlighting}[]
\NormalTok{\textless{}}\KeywordTok{choose}\NormalTok{\textgreater{}}
\NormalTok{  \textless{}}\KeywordTok{if}\OtherTok{ variable=}\StringTok{"issued"}\NormalTok{\textgreater{}}
\NormalTok{    \textless{}}\KeywordTok{date}\OtherTok{ variable=}\StringTok{"issued"}\OtherTok{ form=}\StringTok{"numeric"}\NormalTok{/\textgreater{}}
\NormalTok{  \textless{}/}\KeywordTok{if}\NormalTok{\textgreater{}}
\NormalTok{  \textless{}}\KeywordTok{else}\NormalTok{\textgreater{}}
\NormalTok{    \textless{}}\KeywordTok{text}\OtherTok{ term=}\StringTok{"no date"}\NormalTok{/\textgreater{}}
\NormalTok{  \textless{}/}\KeywordTok{else}\NormalTok{\textgreater{}}
\NormalTok{\textless{}/}\KeywordTok{choose}\NormalTok{\textgreater{}}
\end{Highlighting}
\end{Shaded}

\texttt{cs:choose} requires a \texttt{cs:if} child element, which may be
followed by one or more \texttt{cs:else-if} child elements, and an
optional closing \texttt{cs:else} child element. The \texttt{cs:if} and
\texttt{cs:else-if} elements may contain any number of
\protect\hyperlink{rendering-elements}{rendering elements} (except for
\texttt{cs:layout}). As an empty \texttt{cs:else} element would be
superfluous, \texttt{cs:else} must contain at least one rendering
element. \texttt{cs:if} and \texttt{cs:else-if} elements must carry one
or more conditions, which are set with the attributes:

\begin{description}
\item[\texttt{disambiguate}]
When set to "true" (the only allowed value), the element content is only
rendered if it disambiguates two otherwise identical citations. This
attempt at \protect\hyperlink{disambiguation}{disambiguation} is only
made when all other disambiguation methods have failed to uniquely
identify the target source.
\item[\texttt{is-numeric}]
Tests whether the given variables
(\protect\hyperlink{appendix-iv---variables}{Appendix IV - Variables})
contain numeric content. Content is considered numeric if it solely
consists of numbers. Numbers may have prefixes and suffixes ("D2", "2b",
"L2d"), and may be separated by a comma, hyphen, or ampersand, with or
without spaces ("2, 3", "2-4", "2 \& 4"). For example, "2nd" tests
"true" whereas "second" and "2nd edition" test "false".
\item[\texttt{is-uncertain-date}]
Tests whether the given \protect\hyperlink{date-variables}{date
variables} contain \protect\hyperlink{approximate-dates}{approximate
dates}.
\item[\texttt{locator}]
Tests whether the locator matches the given locator types (see
\protect\hyperlink{locators}{Locators}). Use "sub-verbo" to test for the
"sub verbo" locator type.
\item[\texttt{position}]
Tests whether the cite position matches the given positions
(terminology: citations consist of one or more cites to individual
items). When called within the scope of cs:bibliography,
\texttt{position} tests "false". The positions that can be tested are:

\begin{itemize}
\item
  "first": position of cites that are the first to reference an item
\item
  "ibid"/"ibid-with-locator"/"subsequent": cites referencing previously
  cited items have the "subsequent" position. Such cites may also have
  the "ibid" or "ibid-with-locator" position when:

  \begin{enumerate}
  \def\labelenumi{\alph{enumi})}
  \tightlist
  \item
    the current cite immediately follows on another cite, within the
    same citation, that references the same item
  \end{enumerate}

  or

  \begin{enumerate}
  \def\labelenumi{\alph{enumi})}
  \setcounter{enumi}{1}
  \tightlist
  \item
    the current cite is the first cite in the citation, and the previous
    citation consists of a single cite referencing the same item
  \end{enumerate}

  If either requirement is met, the presence of locators determines
  which position is assigned:

  \begin{itemize}
  \tightlist
  \item
    \textbf{Preceding cite does not have a locator}: if the current cite
    has a locator, the position of the current cite is
    "ibid-with-locator". Otherwise the position is "ibid".
  \item
    \textbf{Preceding cite does have a locator}: if the current cite has
    the same locator, the position of the current cite is "ibid". If the
    locator differs the position is "ibid-with-locator". If the current
    cite lacks a locator its only position is "subsequent".
  \end{itemize}
\item
  "near-note": position of a cite following another cite referencing the
  same item. Both cites have to be located in foot or endnotes, and the
  distance between both cites may not exceed the maximum distance
  (measured in number of foot or endnotes) set with the
  \texttt{near-note-distance} option (see
  \protect\hyperlink{note-distance}{Note Distance}).
\end{itemize}

Whenever position="ibid-with-locator" tests true, position="ibid" also
tests true. And whenever position="ibid" or position="near-note" test
true, position="subsequent" also tests true.
\item[\texttt{type}]
Tests whether the item matches the given types
(\protect\hyperlink{appendix-iii---types}{Appendix III - Types}).
\item[\texttt{variable}]
Tests whether the default (long) forms of the given variables
(\protect\hyperlink{appendix-iv---variables}{Appendix IV - Variables})
contain non-empty values.
\end{description}

With the exception of \texttt{disambiguate}, all conditions allow for
multiple test values (separated with spaces, e.g. "book thesis").

The \texttt{cs:if} and \texttt{cs:else-if} elements may carry the
\texttt{match} attribute to control the testing logic, with allowed
values:

\begin{itemize}
\tightlist
\item
  "all" - (default), element only tests "true" when all conditions test
  "true" for all given test values
\item
  "any" - element tests "true" when any condition tests "true" for any
  given test value
\item
  "none" - element only tests "true" when none of the conditions test
  "true" for any given test value
\end{itemize}

Delimiters from the nearest delimiters from the nearest ancestor
delimiting element \emph{are} applied within the output of
\texttt{cs:choose} (i.e., the output of the matching \texttt{cs:if},
\texttt{cs:else-if}, or \texttt{cs:else}; see
\protect\hyperlink{delimiter}{delimiter}).

\hypertarget{style-behavior}{%
\subsection{Style Behavior}\label{style-behavior}}

\hypertarget{options}{%
\subsubsection{Options}\label{options}}

Styles may be configured with
\protect\hyperlink{citation-specific-options}{citation-specific
options}, set as attributes on set on \texttt{cs:citation},
\protect\hyperlink{bibliography-specific-options}{bibliography-specific
options}, set on \texttt{cs:bibliography}, and
\protect\hyperlink{global-options}{global options} (these affect both
citations and the bibliography), set on \texttt{cs:style}.
\protect\hyperlink{inheritable-name-options}{Inheritable name options}
may be set on \texttt{cs:style}, \texttt{cs:citation} and
\texttt{cs:bibliography}. Finally,
\protect\hyperlink{locale-options}{locale options} may be set on
\texttt{cs:locale} elements.

\hypertarget{citation-specific-options}{%
\paragraph{Citation-specific Options}\label{citation-specific-options}}

\hypertarget{disambiguation}{%
\subparagraph{Disambiguation}\label{disambiguation}}

A cite is ambiguous when it matches multiple bibliographic
entries\footnote{Including uncited entries in the bibliography can make
  cites in the document ambiguous. To make sure such cites are
  disambiguated, the CSL processor should include (invisible) cites for
  such uncited bibliographic entries in the disambiguation process.}.
Four methods are available to eliminate such ambiguity, which are always
tried in the following order:

\begin{enumerate}
\def\labelenumi{\arabic{enumi}.}
\tightlist
\item
  Expand names (adding initials or full given names)
\item
  Show more names
\item
  Render the cite with the \texttt{disambiguate} attribute of
  \texttt{cs:choose} conditions testing "true"
\item
  Add a year-suffix
\end{enumerate}

Method 1 can also be used for the separate purpose of global \emph{name
disambiguation}, covering both ambiguous and unambiguous cites
throughout the document.

The four disambiguation methods can be individually activated with the
following optional attributes:

\begin{description}
\item[\texttt{disambiguate-add-givenname} {[}Method (1){]}]
If set to "true" ("false" is the default), ambiguous names (names that
are identical in their "short" or initialized "long" form, but differ
when initials are added or the full given name is shown) are expanded.
Name expansion can be configured with
\texttt{givenname-disambiguation-rule}. An example of cite
disambiguation:

\begin{longtable}[]{@{}ll@{}}
\toprule\noalign{}
Original ambiguous cites & Disambiguated cites \\
\midrule\noalign{}
\endhead
\bottomrule\noalign{}
\endlastfoot
(Simpson 2005; Simpson 2005) & (H. Simpson 2005; B. Simpson 2005) \\
(Doe 1950; Doe 1950) & (John Doe 1950; Jane Doe 1950) \\
\end{longtable}
\item[\texttt{givenname-disambiguation-rule}]
Specifies (a) whether the purpose of name expansion is limited to
disambiguating cites, or has the additional goal of disambiguating names
(only in the latter case are ambiguous names in unambiguous cites
expanded, e.g. from "(Doe 1950; Doe 2000)" to "(Jane Doe 1950; John Doe
2000)"), (b) whether name expansion targets all, or just the first name
of each cite, and (c) the method by which each name is expanded.

\begin{description}
\item[\textbf{Expansion of Individual Names}]
The steps for expanding individual names are:

\begin{enumerate}
\def\labelenumi{\arabic{enumi}.}
\item
  If \texttt{initialize-with} is set and \texttt{initialize} has its
  default value of "true", then:

  (a) Initials can be shown by rendering the name with a \texttt{form}
  value of "long" instead of "short" (e.g. "Doe" becomes "J. Doe").

  (b) Full given names can be shown instead of initials by rendering the
  name with \texttt{initialize} set to "false" (e.g. "J. Doe" becomes
  "John Doe").
\item
  If \texttt{initialize-with} is \emph{not} set, full given names can be
  shown by rendering the name with a \texttt{form} value of "long"
  instead of "short" (e.g. "Doe" becomes "John Doe").
\end{enumerate}
\item[\textbf{Given Name Disambiguation Rules}]
Allowed values of \texttt{givenname-disambiguation-rule}:

\begin{description}
\item["all-names"]
Name expansion has the dual purpose of disambiguating cites and names.
All rendered ambiguous names, in both ambiguous and unambiguous cites,
are subject to disambiguation. Each name is progressively transformed
until it is disambiguated. Names that cannot be disambiguated remain in
their original form.
\item["all-names-with-initials"]
As "all-names", but name expansion is limited to showing initials (see
step 1(a) above). No disambiguation attempt is made when
\texttt{initialize-with} is not set or when \texttt{initialize} is set
to "false".
\item["primary-name"]
As "all-names", but disambiguation is limited to the first name of each
cite.
\item["primary-name-with-initials"]
As "all-names-with-initials", but disambiguation is limited to the first
name of each cite.
\item["by-cite"]
Default. As "all-names", but the goal of name expansion is limited to
disambiguating cites. Only ambiguous names in ambiguous cites are
affected, and disambiguation stops after the first name that eliminates
cite ambiguity.
\end{description}
\end{description}
\item[\texttt{disambiguate-add-names} {[}Method (2){]}]
If set to "true" ("false" is the default), names that would otherwise be
hidden as a result of et-al abbreviation are added one by one to all
members of a set of ambiguous cites, until no more cites in the set can
be disambiguated by adding names.
\end{description}

If both \texttt{disambiguate-add-givenname} and
\texttt{disambiguate-add-names} are set to "true", given name expansion
is applied to rendered names first. If cites cannot be (fully)
disambiguated by expanding the rendered names, then the names still
hidden as a result of et-al abbreviation are added one by one to all
members of a set of ambiguous cites. Added names are expanded if doing
so would disambiguate the ambiguous cites. This process contines until
no more cites in the set can be disambiguated by adding expanded names.

In the description of disambiguation methods (1) and (2) above, we
assumed that each (disambiguated) cite has an unambiguous link to its
bibliographic entry. To assure that each cite does in fact uniquely
identify its entry in the bibliography, detail that distinguishes cites
(such as names, initials, and full given names) must be shown in the
corresponding bibliography entries. If this is not the case,
disambiguation methods (1) and (2) also act on all members of a set of
ambiguously cited bibliographic entries, until no more entries in the
set can be unambiguously cited by adding (expanded) names. Each method
only takes effect on the involved bibliographic entries after it has
been used to disambiguate cites.

\begin{description}
\item[\texttt{disambiguate} condition {[}Method (3){]}]
A disambiguation attempt can also be made by rendering ambiguous cites
with the \texttt{disambiguate} condition testing "true" (see
\protect\hyperlink{choose}{Choose}).
\item[\texttt{disambiguate-add-year-suffix} {[}Method (4){]}]
If set to "true" ("false" is the default), an alphabetic year-suffix is
added to ambiguous cites (e.g. "Doe 2007, Doe 2007" becomes "Doe 2007a,
Doe 2007b") and to their corresponding bibliographic entries. This final
disambiguation method is always successful. The assignment of
year-suffixes follows the order of the bibliographies entries, and
additional letters are used once "z" is reached ("z", "aa", "ab", ...,
"az", "ba", etc.). By default, the year-suffix is appended the first
year rendered through \texttt{cs:date} in the cite and in the
bibliographic entry, but its location can be controlled by explicitly
rendering the "year-suffix" variable using \texttt{cs:text}. If
"year-suffix" is rendered through \texttt{cs:text} in the scope of
\texttt{cs:citation}, it is suppressed for \texttt{cs:bibliography},
unless it is also rendered through \texttt{cs:text} in the scope of
\texttt{cs:bibliography}, and vice versa.
\end{description}

\hypertarget{cite-grouping}{%
\subparagraph{Cite Grouping}\label{cite-grouping}}

With cite grouping, cites in in-text citations with identical rendered
names are grouped together, e.g. the year-sorted "(Doe 1999; Smith 2002;
Doe 2006; Doe et al. 2007)" becomes "(Doe 1999; Doe 2006; Smith 2002;
Doe et al. 2007)". The comparison is limited to the output of the
(first) \texttt{cs:names} element, but includes output rendered through
\texttt{cs:substitute}. Cite grouping takes places after cite sorting
and disambiguation. Grouped cites maintain their relative order, and are
moved to the original location of the first cite of the group.

Cite grouping can be activated by setting the
\texttt{cite-group-delimiter} attribute or the \texttt{collapse}
attributes on \texttt{cs:citation} (see also
\protect\hyperlink{cite-collapsing}{Cite Collapsing}).

\begin{description}
\item[\texttt{cite-group-delimiter}]
Activates cite grouping and specifies the delimiter for cites within a
cite group. Defaults to ", ". E.g. with \texttt{delimiter} on
\texttt{cs:layout} in \texttt{cs:citation} set to "; ",
\texttt{collapse} set to "year", and \texttt{cite-group-delimiter} set
to ",", citations look like "(Doe 1999,2001; Jones 2000)".
\end{description}

\hypertarget{cite-collapsing}{%
\subparagraph{Cite Collapsing}\label{cite-collapsing}}

Cite groups (author and author-date styles), and numeric cite ranges
(numeric styles) can be collapsed through the use of the
\texttt{collapse} attribute. Delimiters for collapsed cite groups can be
customized with the \texttt{year-suffix-delimiter} and
\texttt{after-collapse-delimiter} attributes:

\begin{description}
\item[\texttt{collapse}]
Activates cite grouping and collapsing. Allowed values:

\begin{itemize}
\tightlist
\item
  "citation-number" - collapses ranges of cite numbers (rendered through
  the "citation-number" variable) in citations for "numeric" styles
  (e.g. from "{[}1, 2, 3, 5{]}" to "{[}1 -- 3, 5{]}"). Only increasing
  ranges collapse, e.g. "{[}3, 2, 1{]}" will not collapse (to see how to
  sort cites by "citation-number", see
  \protect\hyperlink{sorting}{Sorting}).
\item
  "year" - collapses cite groups by suppressing the output of the
  \texttt{cs:names} element for subsequent cites in the group, e.g.
  "(Doe 2000, Doe 2001)" becomes "(Doe 2000, 2001)".
\item
  "year-suffix" - collapses as "year", but also suppresses repeating
  years within the cite group, e.g. "(Doe 2000a, b)" instead of "(Doe
  2000a, 2000b)".
\item
  "year-suffix-ranged" - collapses as "year-suffix", but also collapses
  ranges of year-suffixes, e.g. "(Doe 2000a -- c,e)" instead of "(Doe
  2000a, b, c, e)".
\end{itemize}

"year-suffix" and "year-suffix-ranged" fall back to "year" when
\texttt{disambiguate-add-year-suffix} is "false" (see
\protect\hyperlink{disambiguation}{Disambiguation}), or when a cite has
a locator (e.g. "(Doe 2000a-c, 2000d, p. 5, 2000e,f)", where the cite
for "Doe 2000d" has a locator that prevents the cite from further
collapsing).
\item[\texttt{year-suffix-delimiter}]
Specifies the delimiter for year-suffixes. Defaults to the delimiter set
on \texttt{cs:layout} in \texttt{cs:citation}. E.g. with
\texttt{collapse} set to "year-suffix", \texttt{delimiter} on
\texttt{cs:layout} in \texttt{cs:citation} set to "; ", and
\texttt{year-suffix-delimiter} set to ",", citations look like "(Doe
1999a,b; Jones 2000)".
\item[\texttt{after-collapse-delimiter}]
Specifies the cite delimiter to be used \emph{after} a collapsed cite
group. Defaults to the delimiter set on \texttt{cs:layout} in
\texttt{cs:citation}. E.g. with \texttt{collapse} set to "year",
\texttt{delimiter} on \texttt{cs:layout} in \texttt{cs:citation} set to
", ", and \texttt{after-collapse-delimiter} set to "; ", citations look
like "(Doe 1999, 2001; Jones 2000, Brown 2001)".
\end{description}

\hypertarget{note-distance}{%
\subparagraph{Note Distance}\label{note-distance}}

\begin{description}
\item[\texttt{near-note-distance}]
A cite tests true for the "near-note" position (see
\protect\hyperlink{choose}{Choose}) when a preceding note exists that a)
refers to the same item and b) does not precede the current note by more
footnotes or endnotes than the value of \texttt{near-note-distance}
(default value is "5").
\end{description}

\hypertarget{bibliography-specific-options}{%
\paragraph{Bibliography-specific
Options}\label{bibliography-specific-options}}

\hypertarget{whitespace}{%
\subparagraph{Whitespace}\label{whitespace}}

\begin{description}
\item[\texttt{hanging-indent}]
If set to "true" ("false" is the default), bibliographic entries are
rendered with hanging-indents.
\item[\texttt{second-field-align}]
If set, subsequent lines of bibliographic entries are aligned along the
second field. With "flush", the first field is flush with the margin.
With "margin", the first field is put in the margin, and subsequent
lines are aligned with the margin. An example, where the first field is
\texttt{\textless{}text\ variable="citation-number"\ suffix=".\ "/\textgreater{}}:

\begin{verbatim}
9.  Adams, D. (2002). The Ultimate Hitchhiker's Guide to the
    Galaxy (1st ed.).
10. Asimov, I. (1951). Foundation.
\end{verbatim}
\item[\texttt{line-spacing}]
Specifies vertical line distance. Defaults to "1" (single-spacing), and
can be set to any positive integer to specify a multiple of the standard
unit of line height (e.g. "2" for double-spacing).
\item[\texttt{entry-spacing}]
Specifies vertical distance between bibliographic entries. By default
(with a value of "1"), entries are separated by a single additional
line-height (as set by the line-spacing attribute). Can be set to any
non-negative integer to specify a multiple of this amount.
\end{description}

\hypertarget{reference-grouping}{%
\subparagraph{Reference Grouping}\label{reference-grouping}}

\begin{description}
\item[\texttt{subsequent-author-substitute}]
If set, the value of this attribute replaces names in a bibliographic
entry that also occur in the preceding entry. The exact method of
substitution depends on the value of the
\texttt{subsequent-author-substitute-rule} attribute. Substitution is
limited to the names of the first \texttt{cs:names} element rendered.
\item[\texttt{subsequent-author-substitute-rule}]
Specifies when and how names are substituted as a result of
\texttt{subsequent-author-substitute}. Allowed values:

\begin{itemize}
\tightlist
\item
  "complete-all" - (default), when all names of the name variable match
  those in the preceding bibliographic entry, the value of
  \texttt{subsequent-author-substitute} replaces the entire name list
  (including punctuation and terms like "et al" and "and"), except for
  the affixes set on the \texttt{cs:names} element.
\item
  "complete-each" - requires a complete match like "complete-all", but
  now the value of \texttt{subsequent-author-substitute} substitutes for
  each rendered name.
\item
  "partial-each" - when one or more rendered names in the name variable
  match those in the preceding bibliographic entry, the value of
  \texttt{subsequent-author-substitute} substitutes for each matching
  name. Matching starts with the first name, and continues up to the
  first mismatch.
\item
  "partial-first" - as "partial-each", but substitution is limited to
  the first name of the name variable.
\end{itemize}

For example, take the following bibliographic entries:

\begin{verbatim}
Doe. 1999.
Doe. 2000.
Doe, Johnson & Williams. 2001.
Doe & Smith. 2002.
Doe, Stevens & Miller. 2003.
Doe, Stevens & Miller. 2004.
Doe, Williams et al. 2005.
Doe, Williams et al. 2006.
\end{verbatim}

With \texttt{subsequent-author-substitute} set to "-\/-\/-", and
\texttt{subsequent-author-substitute-rule} set to "complete-all", this
becomes:

\begin{verbatim}
Doe. 1999.
---. 2000.
Doe, Johnson & Williams. 2001.
Doe & Smith. 2002.
Doe, Stevens & Miller. 2003.
---. 2004.
Doe, Williams et al. 2005.
---. 2005.
\end{verbatim}

With \texttt{subsequent-author-substitute-rule} set to "complete-each",
this becomes:

\begin{verbatim}
Doe. 1999.
---. 2000.
Doe, Johnson & Williams. 2001.
Doe & Smith. 2002.
Doe, Stevens & Miller. 2003.
---, --- & ---. 2004.
Doe, Williams et al. 2005.
---, --- et al. 2006.
\end{verbatim}

With \texttt{subsequent-author-substitute-rule} set to "partial-each",
this becomes:

\begin{verbatim}
Doe. 1999.
---. 2000.
Doe, Johnson & Williams. 2001.
--- & Smith. 2002.
Doe, Stevens & Miller. 2003.
---, --- & ---. 2004.
Doe, Williams et al. 2005.
---, --- et al. 2005.
\end{verbatim}

With \texttt{subsequent-author-substitute-rule} set to "partial-first",
this becomes:

\begin{verbatim}
Doe. 1999.
---. 2000.
Doe, Johnson & Williams. 2001.
--- & Smith. 2002.
Doe, Stevens & Miller. 2003.
---, Stevens & Miller. 2004.
Doe, Williams et al. 2005.
---, Williams et al. 2005.
\end{verbatim}
\end{description}

\hypertarget{global-options}{%
\paragraph{Global Options}\label{global-options}}

\hypertarget{hyphenation-of-initialized-names}{%
\subparagraph{Hyphenation of Initialized
Names}\label{hyphenation-of-initialized-names}}

\begin{description}
\item[\texttt{initialize-with-hyphen}]
Specifies whether compound given names (e.g. "Jean-Luc") should be
initialized with a hyphen ("J.-L.", value "true", default) or without
("J. L.", value "false").
\end{description}

\hypertarget{page-ranges}{%
\subparagraph{Page Ranges}\label{page-ranges}}

\begin{description}
\item[\texttt{page-range-format}]
Activates expansion or collapsing of page ranges: "chicago" ("321 --
28"), "expanded" (e.g. "321 -- 328"), "minimal" ("321 -- 8"), or
"minimal-two" ("321 -- 28") (see also
\protect\hyperlink{appendix-v---page-range-formats}{Appendix V - Page
Range Formats}). Delimits page ranges with the "page-range-delimiter"
term (introduced with CSL 1.0.1, and defaults to an en-dash). If the
attribute is not set, page ranges are rendered without reformatting.
\end{description}

\hypertarget{name-particles}{%
\subparagraph{Name Particles}\label{name-particles}}

Western names frequently contain one or more name particles (e.g. "van"
in the Dutch name "Vincent van Gogh"). These name particles can be
either kept or dropped when only the surname is shown: these two types
are referred to as non-dropping and dropping particles, respectively.
Theoretically, a single name might contain particles of both types (with
non-dropping particles always following dropping particles), though
currently we are not aware of any real-life examples. For example, the
Dutch name "Vincent van Gogh", the German name "Alexander von Humboldt",
and the Arabic name "Tawfiq al-Hakim" can be deconstructed into:

\begin{quote}
\begin{verbatim}
{
    "author": [
        {
            "given": "Vincent",
            "non-dropping-particle": "van",
            "family": "Gogh"
        },
        {
            "given": "Alexander",
            "dropping-particle": "von",
            "family": "Humboldt"
        }
        {
            "given": "Tawfiq",
            "non-dropping-particle": "al-",
            "family": "Hakim"
        }
    ]
}
\end{verbatim}
\end{quote}

When just the surname is shown, only the non-dropping-particle is kept:
"Van Gogh" and "al-Hakim", but "Humboldt".

In the case of inverted names, where the family name precedes the given
name, the dropping-particle is always appended to the family name, but
the non-dropping-particle can be either prepended (e.g. "van Gogh,
Vincent") or appended (after initials or given names, e.g. "Gogh,
Vincent van"). For inverted names where the non-dropping-particle is
prepended, names can either be sorted by keeping the
non-dropping-particle together with the family name as part of the
primary sort key (sort order A), or by separating the
non-dropping-particle from the family name and have it become (part of)
a secondary sort key, joining the dropping-particle, if available (sort
order B):

\textbf{Sort order A: non-dropping-particle not demoted}

\begin{itemize}
\tightlist
\item
  primary sort key: "van Gogh"
\item
  secondary sort key: ""
\item
  tertiary sort key: "Vincent"
\end{itemize}

\textbf{Sort order B: non-dropping-particle demoted}

\begin{itemize}
\tightlist
\item
  primary sort key: "Gogh"
\item
  secondary sort key: "van"
\item
  tertiary sort key: "Vincent"
\end{itemize}

The handling of the non-dropping-particle can be customized with the
\texttt{demote-non-dropping-particle} option:

\begin{description}
\item[\texttt{demote-non-dropping-particle}]
Sets the display and sorting behavior of the non-dropping-particle in
inverted names (e.g. "Gogh, Vincent van"). Allowed values:

\begin{itemize}
\tightlist
\item
  "never": the non-dropping-particle is treated as part of the family
  name, whereas the dropping-particle is appended (e.g. "van Gogh,
  Vincent", "Humboldt, Alexander von"). The non-dropping-particle is
  part of the primary sort key (sort order A, e.g. "van Gogh, Vincent"
  appears under "V").
\item
  "sort-only": same display behavior as "never", but the
  non-dropping-particle is demoted to a secondary sort key (sort order
  B, e.g. "van Gogh, Vincent" appears under "G").
\item
  "display-and-sort" (default): the dropping and non-dropping-particle
  are appended (e.g. "Gogh, Vincent van" and "Humboldt, Alexander von").
  For name sorting, all particles are part of the secondary sort key
  (sort order B, e.g. "Gogh, Vincent van" appears under "G").
\end{itemize}
\end{description}

Some names include a particle that should never be demoted. For these
cases the particle should just be included in the family name field, for
example for the French general Charles de Gaulle and the writer Jean de
La Fontaine:

\begin{quote}
\begin{verbatim}
{
    "author": [
        {
            "given": "Charles"
            "family": "de Gaulle",
        },
        {
            "given": "Jean"
            "dropping-particle": "de",
            "family": "La Fontaine",
        }
    ]
}
\end{verbatim}
\end{quote}

\hypertarget{inheritable-name-options}{%
\paragraph{Inheritable Name Options}\label{inheritable-name-options}}

Attributes for the \texttt{cs:names} and \texttt{cs:name} elements may
also be set on \texttt{cs:style}, \texttt{cs:citation} and
\texttt{cs:bibliography}. This eliminates the need to repeat the same
attributes and attribute values for every occurrence of the
\texttt{cs:names} and \texttt{cs:name} elements.

The available inheritable attributes for \texttt{cs:name} are
\texttt{and}, \texttt{delimiter-precedes-et-al},
\texttt{delimiter-precedes-last}, \texttt{et-al-min},
\texttt{et-al-use-first}, \texttt{et-al-use-last},
\texttt{et-al-subsequent-min}, \texttt{et-al-subsequent-use-first},
\texttt{initialize}, \texttt{initialize-with},
\texttt{name-as-sort-order} and \texttt{sort-separator}. The attributes
\texttt{name-form} and \texttt{name-delimiter} correspond to the
\texttt{form} and \texttt{delimiter} attributes on \texttt{cs:name}.
Similarly, \texttt{names-delimiter} corresponds to the
\texttt{delimiter} attribute on \texttt{cs:names}.

When an inheritable name attribute is set on \texttt{cs:style},
\texttt{cs:citation} or \texttt{cs:bibliography}, its value is used for
all \texttt{cs:names} elements within the scope of the element carrying
the attribute. If an attribute is set on multiple hierarchical levels,
the value set at the lowest level is used.

\hypertarget{locale-options}{%
\paragraph{Locale Options}\label{locale-options}}

\begin{description}
\item[\texttt{limit-day-ordinals-to-day-1}]
Date formats are defined by the \texttt{cs:date} element and its
\texttt{cs:date-part} child elements (see
\protect\hyperlink{date}{Date}). By default, when the
\texttt{cs:date-part} element with \texttt{name} set to "day" has
\texttt{form} set to "ordinal", all days (1 through 31) are rendered in
the ordinal form, e.g. "January 1st", "January 2nd", etc. By setting
\texttt{limit-day-ordinals-to-day-1} to "true" ("false" is the default),
the "ordinal" form is limited to the first day of each month (other days
will use the "numeric" form). This is desirable for some languages, such
as French: "1er janvier", but "2 janvier", "3 janvier", etc.
\item[\texttt{punctuation-in-quote}]
For \texttt{cs:text} elements rendered with the \texttt{quotes}
attribute set to "true" (see
\protect\hyperlink{formatting}{Formatting}), and for which the output is
followed by a comma or period, \texttt{punctuation-in-quote} specifies
whether this punctuation is placed outside (value "false", default) or
inside (value "true") the closing quotation mark.
\end{description}

\hypertarget{sorting}{%
\subsubsection{Sorting}\label{sorting}}

\texttt{cs:citation} and \texttt{cs:bibliography} may include a
\texttt{cs:sort} child element before the \texttt{cs:layout} element to
specify the sorting order of respectively cites within citations, and
bibliographic entries within the bibliography. In the absence of
\texttt{cs:sort}, cites and bibliographic entries appear in the order in
which they are cited.

The \texttt{cs:sort} element must contain one or more \texttt{cs:key}
child elements. The sort key, set as an attribute on \texttt{cs:key},
must be a variable (see
\protect\hyperlink{appendix-iv---variables}{Appendix IV - Variables}) or
macro name. For each \texttt{cs:key} element, the sort direction can be
set to either "ascending" (default) or "descending" with the
\texttt{sort} attribute. Sorting is case-insensitive. The attributes
\texttt{names-min}, \texttt{names-use-first}, and
\texttt{names-use-last} may be used to override the values of the
corresponding \texttt{et-al-min}/\texttt{et-al-subsequent-min},
\texttt{et-al-use-first}/\texttt{et-al-subsequent-use-first} and
\texttt{et-al-use-last} attributes, and affect all names generated via
macros called by \texttt{cs:key}.

Sort keys are evaluated in sequence. A primary sort is performed on all
items using the first sort key. A secondary sort, using the second sort
key, is applied to items sharing the first sort key value. A tertiary
sort, using the third sort key, is applied to items sharing the first
and second sort key values. Sorting continues until either the order of
all items is fixed, or until the sort keys are exhausted. Items with an
empty sort key value are placed at the end of the sort, both for
ascending and descending sorts.

An example, where cites are first sorted by the output of the "author"
macro, with overriding settings for et-al abbreviation. Cites sharing
the primary sort key are subsequently sorted in descending order by the
"issued" date variable.

\begin{Shaded}
\begin{Highlighting}[]
\NormalTok{\textless{}}\KeywordTok{citation}\NormalTok{\textgreater{}}
\NormalTok{  \textless{}}\KeywordTok{sort}\NormalTok{\textgreater{}}
\NormalTok{    \textless{}}\KeywordTok{key}\OtherTok{ macro=}\StringTok{"author"}\OtherTok{ names{-}min=}\StringTok{"3"}\OtherTok{ names{-}use{-}first=}\StringTok{"3"}\NormalTok{/\textgreater{}}
\NormalTok{    \textless{}}\KeywordTok{key}\OtherTok{ variable=}\StringTok{"issued"}\OtherTok{ sort=}\StringTok{"descending"}\NormalTok{/\textgreater{}}
\NormalTok{  \textless{}/}\KeywordTok{sort}\NormalTok{\textgreater{}}
\NormalTok{  \textless{}}\KeywordTok{layout}\NormalTok{\textgreater{}}
    \CommentTok{\textless{}!{-}{-} rendering elements {-}{-}\textgreater{}}
\NormalTok{  \textless{}/}\KeywordTok{layout}\NormalTok{\textgreater{}}
\NormalTok{\textless{}/}\KeywordTok{citation}\NormalTok{\textgreater{}}
\end{Highlighting}
\end{Shaded}

The sort key value of a variable or macro can differ from the "normal"
rendered output. The specifics of sorting variables and macros:

\hypertarget{sorting-variables}{%
\paragraph{Sorting Variables}\label{sorting-variables}}

The sort key value for a variable called by \texttt{cs:key} via the
\texttt{variable} attribute consists of the string value, without rich
text markup. Exceptions are name, date and numeric variables:

\textbf{names:} \protect\hyperlink{name-variables}{Name variables}
called via the \texttt{variable} attribute (e.g.
\texttt{\textless{}key\ variable="author"/\textgreater{}}) are returned
as a name list string, with the \texttt{cs:name} attributes
\texttt{form} set to "long", and \texttt{name-as-sort-order} set to
"all".

\textbf{dates:} \protect\hyperlink{date-variables}{Date variables}
called via the \texttt{variable} attribute are returned in the YYYYMMDD
format, with zeros substituted for any missing date-parts (e.g. 20001200
for December 2000). As a result, less specific dates precede more
specific dates in ascending sorts, e.g. "2000, May 2000, May 1st 2000".
Negative years are sorted inversely, e.g. "100BC, 50BC, 50AD, 100AD".
Seasons are ignored for sorting, as the chronological order of the
seasons differs between the northern and southern hemispheres. In the
case of date ranges, the start date is used for the primary sort, and
the end date is used for a secondary sort, e.g. "2000 -- 2001, 2000 --
2005, 2002 -- 2003, 2002 -- 2009". Date ranges are placed after single
dates when they share the same (start) date, e.g. "2000, 2000 -- 2002".

\textbf{numbers:} \protect\hyperlink{number-variables}{Number variables}
called via the \texttt{variable} attribute are returned as integers
(\texttt{form} is "numeric"). If the original variable value only
consists of non-numeric text, the value is returned as a text string.

\hypertarget{sorting-macros}{%
\paragraph{Sorting Macros}\label{sorting-macros}}

The sort key value for a macro called via \texttt{cs:key} via the
\texttt{macro} attribute generally consists of the string value the
macro would ordinarily generate, without rich text markup (exceptions
are discussed below).

When sorting by \protect\hyperlink{name-variables}{name variables}
rendered within the macro, any \texttt{cs:label} elements are excluded
from the sort key values. Name variables are also returned with the
\texttt{cs:name} attribute \texttt{name-as-sort-order} set to "all".
When et-al abbreviation occurs, the "et-al" and "and others" terms are
excluded from the sort key values.

There are four advantages in using the same macro for rendering and
sorting, instead of sorting directly on the name variable. First,
substitution is available (e.g. the "editor" variable might substitute
for an empty "author" variable). Second, et-al abbreviation can be used
(using either the \texttt{et-al-min}/\texttt{et-al-subsequent-min},
\texttt{et-al-use-first}/\texttt{et-al-subsequent-use-first}, and
\texttt{et-al-use-last} options defined for the macro, or the overriding
\texttt{names-min}, \texttt{names-use-first} and \texttt{names-use-last}
attributes set on \texttt{cs:key}). Third, names can be sorted by just
the surname (using a macro for which the \texttt{form} attribute on
\texttt{cs:name} is set to "short"). Fourth, it is possible to sort by
the number of names in a name list, by calling a macro for which the
\texttt{form} attribute on \texttt{cs:name} is set to "count".

\protect\hyperlink{number-variables}{Number variables} rendered within
the macro with \texttt{cs:number} and
\protect\hyperlink{date-variables}{date variables} are treated the same
as when they are called via \texttt{variable}. The only exception is
that the complete date is returned if a date variable is called via the
\texttt{variable} attribute. In contrast, macros return only those
date-parts that would otherwise be rendered (respecting the value of the
\texttt{date-parts} attribute for localized dates, or the listing of
\texttt{cs:date-part} elements for non-localized dates).

\hypertarget{range-delimiters}{%
\subsubsection{Range Delimiters}\label{range-delimiters}}

Collapsed ranges for the "citation-number" and "year-suffix" variables
are delimited by an en-dash (e.g. "(1 -- 3, 5)" and "(Doe 2000a --
c,e)").

The "locator" variable is always rendered with an en-dash replacing any
hyphens. For the "page" variable, this replacement is only performed if
the \texttt{page-range-format} attribute is set on \texttt{cs:style}
(see \protect\hyperlink{page-ranges}{Page Ranges}).

\hypertarget{formatting}{%
\subsubsection{Formatting}\label{formatting}}

The following formatting attributes may be set on \texttt{cs:date},
\texttt{cs:date-part}, \texttt{cs:et-al}, \texttt{cs:group},
\texttt{cs:label}, \texttt{cs:layout}, \texttt{cs:name},
\texttt{cs:name-part}, \texttt{cs:names}, \texttt{cs:number} and
\texttt{cs:text}:

\begin{description}
\item[\texttt{font-style}]
Sets the font style, with values:

\begin{itemize}
\tightlist
\item
  "normal" (default)
\item
  "italic"
\item
  "oblique" (i.e. slanted)
\end{itemize}
\item[\texttt{font-variant}]
Allows for the use of small capitals, with values:

\begin{itemize}
\tightlist
\item
  "normal" (default)
\item
  "small-caps"
\end{itemize}
\item[\texttt{font-weight}]
Sets the font weight, with values:

\begin{itemize}
\tightlist
\item
  "normal" (default)
\item
  "bold"
\item
  "light"
\end{itemize}
\item[\texttt{text-decoration}]
Allows for the use of underlining, with values:

\begin{itemize}
\tightlist
\item
  "none" (default)
\item
  "underline"
\end{itemize}
\item[\texttt{vertical-align}]
Sets the vertical alignment, with values:

\begin{itemize}
\tightlist
\item
  "baseline" (default)
\item
  "sup" (superscript)
\item
  "sub" (subscript)
\end{itemize}
\end{description}

\hypertarget{affixes}{%
\subsubsection{Affixes}\label{affixes}}

The affixes attributes \texttt{prefix} and \texttt{suffix} may be set on
\texttt{cs:date} (except when \texttt{cs:date} defines a localized date
format), \texttt{cs:date-part} (except when the parent \texttt{cs:date}
element calls a localized date format), \texttt{cs:group},
\texttt{cs:label}, \texttt{cs:layout}, \texttt{cs:name},
\texttt{cs:name-part}, \texttt{cs:names}, \texttt{cs:number} and
\texttt{cs:text}. The attribute value is either added before
(\texttt{prefix}) or after (\texttt{suffix}) the output of the element
carrying the attribute, but affixes are only rendered if the element
produces output. With the exception of affixes set on
\texttt{cs:layout}, affixes are outside the scope of
\protect\hyperlink{formatting}{formatting},
\protect\hyperlink{quotes}{quotes},
\protect\hyperlink{strip-periods}{strip-periods} and
\protect\hyperlink{text-case}{text-case} attributes set on the same
element (as a workaround, these attributes take effect on affixes when
set on a parent \texttt{cs:group} element).

\hypertarget{delimiter}{%
\subsubsection{Delimiter}\label{delimiter}}

The \texttt{delimiter} attribute, whose value delimits non-empty pieces
of output, may be set on \texttt{cs:date} (delimiting the date-parts;
\texttt{delimiter} is not allowed when \texttt{cs:date} calls a
localized date format), \texttt{cs:names} (delimiting name lists from
different \protect\hyperlink{name-variables}{name variables}),
\texttt{cs:name} (delimiting names within name lists), \texttt{cs:group}
and \texttt{cs:layout} (delimiting the output of the child elements).

A delimiting element is any element as above which takes a
\texttt{delimiter} attribute, whether the attribute is supplied or not.

Delimiters from any ancestor delimiting element are not applied within
the output of a delimiting element. The following produces
\texttt{retrieved:\ \textless{}http://example.com\textgreater{}}:

\begin{Shaded}
\begin{Highlighting}[]
\NormalTok{\textless{}}\KeywordTok{group}\OtherTok{ delimiter=}\StringTok{": "}\NormalTok{\textgreater{}}
\NormalTok{  \textless{}}\KeywordTok{text}\OtherTok{ term=}\StringTok{"retrieved"}\NormalTok{ /\textgreater{}}
\NormalTok{  \textless{}}\KeywordTok{group}\NormalTok{\textgreater{}}
\NormalTok{    \textless{}}\KeywordTok{text}\OtherTok{ value=}\StringTok{"}\DecValTok{\&lt;}\StringTok{"}\NormalTok{ /\textgreater{}}
\NormalTok{    \textless{}}\KeywordTok{text}\OtherTok{ variable=}\StringTok{"URL"}\NormalTok{ /\textgreater{}}
\NormalTok{    \textless{}}\KeywordTok{text}\OtherTok{ value=}\StringTok{"}\DecValTok{\&gt;}\StringTok{"}\NormalTok{ /\textgreater{}}
\NormalTok{  \textless{}/}\KeywordTok{group}\NormalTok{\textgreater{}}
\NormalTok{\textless{}/}\KeywordTok{group}\NormalTok{\textgreater{}}
\end{Highlighting}
\end{Shaded}

\hypertarget{display}{%
\subsubsection{Display}\label{display}}

The \texttt{display} attribute (similar the "display" property in CSS)
may be used to structure individual bibliographic entries into one or
more text blocks. If used, all rendering elements should be under the
control of a display attribute. The allowed values:

\begin{itemize}
\tightlist
\item
  "block" - block stretching from margin to margin.
\item
  "left-margin" - block starting at the left margin. If followed by a
  "right-inline" block, the "left-margin" blocks of all bibliographic
  entries are set to a fixed width to accommodate the longest content
  string found among these "left-margin" blocks. In the absence of a
  "right-inline" block the "left-margin" block extends to the right
  margin.
\item
  "right-inline" - block starting to the right of a preceding
  "left-margin" block (behaves as "block" in the absence of such a
  "left-margin" block). Extends to the right margin.
\item
  "indent" - block indented to the right by a standard amount. Extends
  to the right margin.
\end{itemize}

\textbf{Examples}

\begin{enumerate}
\def\labelenumi{(\Alph{enumi})}
\item
  Instead of using \texttt{second-field-align} (see
  \protect\hyperlink{whitespace}{Whitespace}), a similar layout can be
  achieved with a "left-margin" and "right-inline" block. A potential
  benefit is that the styling of blocks can be further controlled in the
  final output (e.g. using CSS for HTML, styles for Word, etc.).

\begin{Shaded}
\begin{Highlighting}[]
\NormalTok{\textless{}}\KeywordTok{bibliography}\NormalTok{\textgreater{}}
\NormalTok{  \textless{}}\KeywordTok{layout}\NormalTok{\textgreater{}}
\NormalTok{    \textless{}}\KeywordTok{text}\OtherTok{ display=}\StringTok{"left{-}margin"}\OtherTok{ variable=}\StringTok{"citation{-}number"}
\OtherTok{        prefix=}\StringTok{"["}\OtherTok{ suffix=}\StringTok{"]"}\NormalTok{/\textgreater{}}
\NormalTok{    \textless{}}\KeywordTok{group}\OtherTok{ display=}\StringTok{"right{-}inline"}\NormalTok{\textgreater{}}
      \CommentTok{\textless{}!{-}{-} rendering elements {-}{-}\textgreater{}}
\NormalTok{    \textless{}/}\KeywordTok{group}\NormalTok{\textgreater{}}
\NormalTok{  \textless{}/}\KeywordTok{layout}\NormalTok{\textgreater{}}
\NormalTok{\textless{}/}\KeywordTok{bibliography}\NormalTok{\textgreater{}}
\end{Highlighting}
\end{Shaded}
\end{enumerate}

\begin{center}\rule{0.5\linewidth}{0.5pt}\end{center}

\begin{enumerate}
\def\labelenumi{(\Alph{enumi})}
\setcounter{enumi}{1}
\item
  A per-author publication listing. With
  \texttt{subsequent-author-substitute} (see
  \protect\hyperlink{reference-grouping}{Reference Grouping}) set to an
  empty string, the block with the author names is only rendered once
  for items by the same authors.

\begin{Shaded}
\begin{Highlighting}[]
\NormalTok{\textless{}}\KeywordTok{bibliography}\OtherTok{ subsequent{-}author{-}substitute=}\StringTok{""}\NormalTok{\textgreater{}}
\NormalTok{  \textless{}}\KeywordTok{sort}\NormalTok{\textgreater{}}
\NormalTok{    \textless{}}\KeywordTok{key}\OtherTok{ variable=}\StringTok{"author"}\NormalTok{/\textgreater{}}
\NormalTok{    \textless{}}\KeywordTok{key}\OtherTok{ variable=}\StringTok{"issued"}\NormalTok{/\textgreater{}}
\NormalTok{  \textless{}/}\KeywordTok{sort}\NormalTok{\textgreater{}}
\NormalTok{  \textless{}}\KeywordTok{layout}\NormalTok{\textgreater{}}
\NormalTok{    \textless{}}\KeywordTok{group}\OtherTok{ display=}\StringTok{"block"}\NormalTok{\textgreater{}}
\NormalTok{      \textless{}}\KeywordTok{names}\OtherTok{ variable=}\StringTok{"author"}\NormalTok{/\textgreater{}}
\NormalTok{    \textless{}/}\KeywordTok{group}\NormalTok{\textgreater{}}
\NormalTok{    \textless{}}\KeywordTok{group}\OtherTok{ display=}\StringTok{"left{-}margin"}\NormalTok{\textgreater{}}
\NormalTok{      \textless{}}\KeywordTok{date}\OtherTok{ variable=}\StringTok{"issued"}\NormalTok{\textgreater{}}
\NormalTok{        \textless{}}\KeywordTok{date{-}part}\OtherTok{ name=}\StringTok{"year"}\NormalTok{ /\textgreater{}}
\NormalTok{      \textless{}/}\KeywordTok{date}\NormalTok{\textgreater{}}
\NormalTok{    \textless{}/}\KeywordTok{group}\NormalTok{\textgreater{}}
\NormalTok{    \textless{}}\KeywordTok{group}\OtherTok{ display=}\StringTok{"right{-}inline"}\NormalTok{\textgreater{}}
\NormalTok{      \textless{}}\KeywordTok{text}\OtherTok{ variable=}\StringTok{"title"}\NormalTok{/\textgreater{}}
\NormalTok{    \textless{}/}\KeywordTok{group}\NormalTok{\textgreater{}}
\NormalTok{  \textless{}/}\KeywordTok{layout}\NormalTok{\textgreater{}}
\NormalTok{\textless{}/}\KeywordTok{bibliography}\NormalTok{\textgreater{}}
\end{Highlighting}
\end{Shaded}

  The output of this example would look like:

  \begin{longtable}[]{@{}
    >{\raggedright\arraybackslash}p{(\columnwidth - 2\tabcolsep) * \real{0.2778}}
    >{\raggedright\arraybackslash}p{(\columnwidth - 2\tabcolsep) * \real{0.3333}}@{}}
  \toprule\noalign{}
  \endhead
  \bottomrule\noalign{}
  \endlastfoot
  \multicolumn{2}{@{}>{\raggedright\arraybackslash}p{(\columnwidth - 2\tabcolsep) * \real{0.6111} + 2\tabcolsep}@{}}{%
  Author1} \\
  year-publication1 & title-publication1 \\
  year-publication2 & title-publication2 \\
  \multicolumn{2}{@{}>{\raggedright\arraybackslash}p{(\columnwidth - 2\tabcolsep) * \real{0.6111} + 2\tabcolsep}@{}}{%
  Author2} \\
  year-publication3 & title-publication3 \\
  year-publication4 & title-publication4 \\
  \end{longtable}
\end{enumerate}

\begin{center}\rule{0.5\linewidth}{0.5pt}\end{center}

\begin{enumerate}
\def\labelenumi{(\Alph{enumi})}
\setcounter{enumi}{2}
\item
  An annotated bibliography, where the annotation appears in an indented
  block below the reference.

\begin{Shaded}
\begin{Highlighting}[]
\NormalTok{\textless{}}\KeywordTok{bibliography}\NormalTok{\textgreater{}}
\NormalTok{  \textless{}}\KeywordTok{layout}\NormalTok{\textgreater{}}
\NormalTok{    \textless{}}\KeywordTok{group}\OtherTok{ display=}\StringTok{"block"}\NormalTok{\textgreater{}}
      \CommentTok{\textless{}!{-}{-} rendering elements {-}{-}\textgreater{}}
\NormalTok{    \textless{}/}\KeywordTok{group}\NormalTok{\textgreater{}}
\NormalTok{    \textless{}}\KeywordTok{text}\OtherTok{ display=}\StringTok{"indent"}\OtherTok{ variable=}\StringTok{"abstract"}\NormalTok{ /\textgreater{}}
\NormalTok{  \textless{}/}\KeywordTok{layout}\NormalTok{\textgreater{}}
\NormalTok{\textless{}/}\KeywordTok{bibliography}\NormalTok{\textgreater{}}
\end{Highlighting}
\end{Shaded}
\end{enumerate}

\hypertarget{quotes}{%
\subsubsection{Quotes}\label{quotes}}

The \texttt{quotes} attribute may set on \texttt{cs:text}. When set to
"true" ("false" is default), the rendered text is wrapped in quotes (the
quotation marks used are terms). The localized
\texttt{punctuation-in-quote} option controls whether an adjoining comma
or period appears outside (default) or inside the closing quotation mark
(see \protect\hyperlink{locale-options}{Locale Options}).

\hypertarget{strip-periods}{%
\subsubsection{Strip-periods}\label{strip-periods}}

The \texttt{strip-periods} attribute may be set on \texttt{cs:date-part}
(but only if \texttt{name} is set to "month"), \texttt{cs:label} and
\texttt{cs:text}. When set to "true" ("false" is the default), any
periods in the rendered text are removed.

\hypertarget{text-case}{%
\subsubsection{Text-case}\label{text-case}}

The \texttt{text-case} attribute may be set on \texttt{cs:date},
\texttt{cs:date-part}, \texttt{cs:label}, \texttt{cs:name-part},
\texttt{cs:number} and \texttt{cs:text}. The allowed values:

\begin{itemize}
\tightlist
\item
  "lowercase": renders text in lowercase
\item
  "uppercase": renders text in uppercase
\item
  "capitalize-first": capitalizes the first character of the first word,
  if the word is lowercase
\item
  "capitalize-all": capitalizes the first character of every lowercase
  word
\item
  "sentence": renders text in sentence case (deprecated; do not use)
\item
  "title": renders text in title case
\end{itemize}

\hypertarget{sentence-case-conversion}{%
\paragraph{Sentence Case Conversion}\label{sentence-case-conversion}}

Sentence case conversion (with \texttt{text-case} set to "sentence") is
performed by:

\begin{enumerate}
\def\labelenumi{\arabic{enumi}.}
\tightlist
\item
  For uppercase strings, the first character of the string remains
  capitalized. All other letters are lowercased.
\item
  For lower or mixed case strings, the first character of the first word
  is capitalized if the word is lowercase. The case of all other words
  stays the same.
\end{enumerate}

CSL processors don\textquotesingle t recognize proper nouns. As a
result, strings in sentence case can be accurately converted to title
case, but not vice versa. For this reason, it is generally preferable to
store strings such as titles in sentence case, and only use
\texttt{text-case} if a style desires another case.

Sentence case conversion is deprecated and will be removed in a future
version.

\hypertarget{title-case-conversion}{%
\paragraph{Title Case Conversion}\label{title-case-conversion}}

Title case conversion (with \texttt{text-case} set to "title") for
English-language items is performed by:

\begin{enumerate}
\def\labelenumi{\arabic{enumi}.}
\tightlist
\item
  For uppercase strings, the first character of each word remains
  capitalized. All other letters are lowercased.
\item
  For lower or mixed case strings, the first character of each lowercase
  word is capitalized. The case of words in mixed or uppercase stays the
  same.
\end{enumerate}

In both cases, stop words are lowercased, unless they are the first or
last word in the string, or follow a colon. The stop words are listed in
the CSL Schema file
\href{https://resource.citationstyles.org/schema/latest/styles/stop-words.json}{stop-words.json}.
Hyphenated word parts are treated as distinct words (e.g., "two-thirds"
becomes "Two-Thirds").

\hypertarget{non-english-items}{%
\subparagraph{Non-English Items}\label{non-english-items}}

As many languages do not use title case, title case conversion (with
\texttt{text-case} set to "title") only affects English-language items.

If the \texttt{default-locale} attribute on \texttt{cs:style}
isn\textquotesingle t set, or set to a locale code with a primary
language tag of "en" (English), items are assumed to be English. An item
is only considered to be non-English if its metadata contains a
\texttt{language} field with a non-nil value that
doesn\textquotesingle t start with the "en" primary language tag.

If \texttt{default-locale} is set to a locale code with a primary
language tag other than "en", items are assumed to be non-English. An
item is only considered to be English if the value of its
\texttt{language} field starts with the "en" primary language tag.

\hypertarget{appendix-i---categories}{%
\subsection{Appendix I - Categories}\label{appendix-i---categories}}

\begin{itemize}
\tightlist
\item
  anthropology
\item
  astronomy
\item
  biology
\item
  botany
\item
  chemistry
\item
  communications
\item
  engineering
\item
  generic-base - used for generic styles like Harvard and APA
\item
  geography
\item
  geology
\item
  history
\item
  humanities
\item
  law
\item
  linguistics
\item
  literature
\item
  math
\item
  medicine
\item
  philosophy
\item
  physics
\item
  political\_science
\item
  psychology
\item
  science
\item
  social\_science
\item
  sociology
\item
  theology
\item
  zoology
\end{itemize}

\hypertarget{appendix-ii---terms}{%
\subsection{Appendix II - Terms}\label{appendix-ii---terms}}

\hypertarget{type-terms}{%
\subsubsection{Type Terms}\label{type-terms}}

For each item type listed in
\protect\hyperlink{appendix-iii---types}{Appendix III - Types}, there is
a corresponding term.

\hypertarget{name-variable-terms}{%
\subsubsection{Name Variable Terms}\label{name-variable-terms}}

For each of the \protect\hyperlink{name-variables}{Name Variables}
listed in \protect\hyperlink{appendix-iv---variables}{Appendix IV -
Variables}, there is a corresponding term.

\hypertarget{number-variable-terms}{%
\subsubsection{Number Variable Terms}\label{number-variable-terms}}

For each of the \protect\hyperlink{number-variables}{Number Variables}
listed in \protect\hyperlink{appendix-iv---variables}{Appendix IV -
Variables}, there is a corresponding term.

\hypertarget{locators}{%
\subsubsection{Locators}\label{locators}}

\begin{itemize}
\tightlist
\item
  appendix
\item
  article-locator
\item
  book
\item
  canon
\item
  chapter
\item
  column
\item
  elocation
\item
  equation
\item
  figure
\item
  folio
\item
  issue
\item
  line
\item
  note
\item
  opus
\item
  page
\item
  paragraph
\item
  part
\item
  rule
\item
  section
\item
  sub-verbo
\item
  supplement
\item
  table
\item
  timestamp
\item
  title
\item
  verse
\item
  volume
\end{itemize}

\hypertarget{months}{%
\subsubsection{Months}\label{months}}

\begin{itemize}
\tightlist
\item
  month-01
\item
  month-02
\item
  month-03
\item
  month-04
\item
  month-05
\item
  month-06
\item
  month-07
\item
  month-08
\item
  month-09
\item
  month-10
\item
  month-11
\item
  month-12
\end{itemize}

\hypertarget{ordinals}{%
\subsubsection{Ordinals}\label{ordinals}}

\begin{itemize}
\tightlist
\item
  ordinal
\item
  ordinal-00 through ordinal-99
\item
  long-ordinal-01
\item
  long-ordinal-02
\item
  long-ordinal-03
\item
  long-ordinal-04
\item
  long-ordinal-05
\item
  long-ordinal-06
\item
  long-ordinal-07
\item
  long-ordinal-08
\item
  long-ordinal-09
\item
  long-ordinal-10
\end{itemize}

\hypertarget{punctuation}{%
\subsubsection{Punctuation}\label{punctuation}}

\begin{itemize}
\tightlist
\item
  open-quote
\item
  close-quote
\item
  open-inner-quote
\item
  close-inner-quote
\item
  page-range-delimiter
\item
  colon
\item
  comma
\item
  semicolon
\end{itemize}

\hypertarget{seasons-1}{%
\subsubsection{Seasons}\label{seasons-1}}

\begin{itemize}
\tightlist
\item
  season-01
\item
  season-02
\item
  season-03
\item
  season-04
\end{itemize}

\hypertarget{miscellaneous}{%
\subsubsection{Miscellaneous}\label{miscellaneous}}

\begin{itemize}
\tightlist
\item
  accessed
\item
  ad
\item
  advance-online-publication
\item
  album
\item
  and
\item
  and others
\item
  anonymous
\item
  at
\item
  audio-recording
\item
  available at
\item
  bc
\item
  bce
\item
  by
\item
  ce
\item
  circa
\item
  cited
\item
  et-al
\item
  film
\item
  forthcoming
\item
  from
\item
  henceforth
\item
  ibid
\item
  in
\item
  in press
\item
  internet
\item
  interview
\item
  letter
\item
  loc-cit
\item
  no date
\item
  no-place
\item
  no-publisher
\item
  on
\item
  online
\item
  op-cit
\item
  original-work-published
\item
  personal-communication
\item
  podcast
\item
  podcast-episode
\item
  preprint
\item
  presented at
\item
  radio-broadcast
\item
  radio-series
\item
  radio-series-episode
\item
  reference
\item
  retrieved
\item
  review-of
\item
  scale
\item
  special-issue
\item
  special-section
\item
  television-broadcast
\item
  television-series
\item
  television-series-episode
\item
  video
\item
  working-paper
\end{itemize}

\hypertarget{appendix-iii---types}{%
\subsection{Appendix III - Types}\label{appendix-iii---types}}

\begin{description}

  \item[article]
    A self-contained work made widely available but not published in a
    journal or other publication;\\
    Use for preprints, working papers, and similar works posted on a
    platform where some level of persistence or stewardship is expected
    (e.g. arXiv or other preprint repositories, working paper series);\\
    For unpublished works not made widely available or only hosted on
    personal websites, use \texttt{manuscript}

  \item[article-journal]
    An article published in an academic journal

  \item[article-magazine]
    An article published in a non-academic magazine

  \item[article-newspaper]
    An article published in a newspaper

  \item[bill]
    A proposed piece of legislation

  \item[book]
    A book or similar work;\\
    Can be an authored book or an edited collection of self-contained
    chapters;\\
    Can be a physical book or an ebook;\\
    The format for an ebook may be specified using \texttt{medium};\\
    Can be a single-volume work, a multivolume work, or one volume of a
    multivolume work;\\
    If a \texttt{container-title} is present, the item is interpreted as a
    book republished in a collection or anthology;\\
    Also used for whole conference proceedings volumes or exhibition
    catalogs by specifying \texttt{event} and related variables

  \item[broadcast]
    A recorded work broadcast over an electronic medium (e.g. a radio
    broadcast, a television show, a podcast);\\
    The type of broadcast may be specified using \texttt{genre};\\
    If \texttt{container-title} is present, the item is interpreted as an
    episode contained within a larger broadcast series (e.g. an episode in a
    television show or an episode of a podcast)

  \item[chapter]
    A part of a book cited separately from the book as a whole (e.g. a
    chapter in an edited book);\\
    Also used for introductions, forewords, and similar supplemental
    components of a book

  \item[classic]
    A classical or ancient work, sometimes cited using a common abbreviation

  \item[collection]
    An archival collection in a museum or other institution

  \item[dataset]
    A data set or a similar collection of (mostly) raw data

  \item[document]
    A catch-all category for items not belonging to other types;\\
    Use a more specific type when appropriate

  \item[entry]
    An entry in a database, directory, or catalog;\\
    For entries in a dictionary, use \texttt{entry-dictionary};\\
    For entries in an encyclopedia, use \texttt{entry-encyclopedia}

  \item[entry-dictionary]
    An entry in a dictionary

  \item[entry-encyclopedia]
    An entry in an encyclopedia or similar reference work

  \item[event]
    An organized event (e.g., an exhibition or conference);\\
    Use for direct citations to the event, rather than to works contained
    within an event (e.g. a \texttt{presentation} in a conference, a
    \texttt{graphic} in an exhibition) or based on an event (e.g. a
    \texttt{paper-conference} published in a proceedings, an exhibition
    catalog)

  \item[figure]
    A illustration or representation of data, typically as part of a journal
    article or other larger work;\\
    May be in any format (e.g. image, video, audio recording, 3D model);\\
    The format of the item can be specified using \texttt{medium}

  \item[graphic]
    A still visual work;\\
    Can be used for artwork or other works (e.g. journalistic or historical
    photographs);\\
    Can be used for any still visual work (e.g. photographs, drawings,
    paintings, sculptures, clothing);\\
    The format of the item can be specified using \texttt{medium}

  \item[hearing]
    A hearing by a government committee or transcript thereof

  \item[interview]
    An interview of a person;\\
    Also used for a recording or transcript of an interview; \texttt{author}
    is interpreted as the interviewee

  \item[legal\_case]
    A legal case

  \item[legislation]
    A law or resolution enacted by a governing body

  \item[manuscript]
    An unpublished manuscript;\\
    Use for both modern unpublished works and classical manuscripts;\\
    For working papers, preprints, and similar works posted to a repository,
    use \texttt{article}

  \item[map]
    A geographic map

  \item[motion\_picture]
    A video or visual recording;\\
    If a \texttt{container-title} is present, the item is interpreted as a
    part contained within a larger compilation of recordings (e.g. a part of
    a multipart documentary))

  \item[musical\_score]
    The printed score for a piece of music;\\
    For a live performance of the music, use \texttt{performance};\\
    For recordings of the music, use \texttt{song} (for audio recordings) or
    \texttt{motion\_picture} (for video recordings)

  \item[pamphlet]
    A fragment, historical document, or other unusually-published or
    ephemeral work (e.g. a sales brochure)

  \item[paper-conference]
    A paper formally published in conference proceedings;\\
    For papers presented at a conference, but not published in a
    proceedings, use \texttt{speech}

  \item[patent]
    A patent for an invention

  \item[performance]
    A live performance of an artistic work;\\
    For non-artistic presentations, use \texttt{speech};\\
    For recordings of a performance, use \texttt{song} or
    \texttt{motion\_picture}

  \item[periodical]
    A full issue or run of issues in a periodical publication (e.g. a
    special issue of a journal)

  \item[personal\_communication]
    Personal communications between multiple parties;\\
    May be unpublished (e.g. private correspondence between two researchers)
    or collected/published (e.g. a letter published in a collection)

  \item[post]
    A post on a online forum, social media platform, or similar platform;\\
    Also used for comments posted to online items

  \item[post-weblog]
    A blog post

  \item[regulation]
    An administrative order from any level of government

  \item[report]
    A technical report, government report, white paper, brief, or similar
    work distributed by an institution;\\
    Also used for manuals and similar technical documentation (e.g. a
    software, instrument, or test manual);\\
    If a \texttt{container-title} is present, the item is interpreted as a
    chapter contained within a larger report

  \item[review]
    A review of an item other than a book (e.g. a film review, posted peer
    review of an article);\\
    If \texttt{reviewed-title} is absent, \texttt{title} is taken to be the
    title of the reviewed item

  \item[review-book]
    A review of a book;\\
    If \texttt{reviewed-title} is absent, \texttt{title} is taken to be the
    title of the reviewed book

  \item[software]
    A computer program, app, or other piece of software

  \item[song]
    An audio recording;\\
    Can be used for any audio recording (not only music);\\
    If a \texttt{container-title} is present, the item is interpreted as a
    track contained within a larger album or compilation of recordings

  \item[speech]
    A speech or other presentation (e.g. a paper, talk, poster, or symposium
    at a conference);\\
    Use \texttt{genre} to specify the type of presentation;\\
    Use \texttt{event} to indicate the event where the presentation was made
    (e.g. the conference name);\\
    Use \texttt{container-title} if the presentation is part of a larger
    session (e.g. a paper in a symposium);\\
    For papers published in conference proceedings, use
    \texttt{paper-conference};\\
    For artistic performances, use \texttt{performance}

  \item[standard]
    A technical standard or similar set of rules or norms

  \item[thesis]
    A thesis written to satisfy requirements for a degree;\\
    Use \texttt{genre} to specify the type of thesis

  \item[treaty]
    A treaty agreement among political authorities

  \item[webpage]
    A website or page on a website;\\
    Intended for sources which are intrinsically online; use a more specific
    type when appropriate (e.g. \texttt{article-journal},
    \texttt{post-weblog}, \texttt{report}, \texttt{entry});\\
    If a \texttt{container-title} is present, the item is interpreted as a
    page contained within a larger website

\end{description}

\hypertarget{appendix-iv---variables}{%
\subsection{Appendix IV - Variables}\label{appendix-iv---variables}}

\hypertarget{standard-variables}{%
\subsubsection{Standard Variables}\label{standard-variables}}

\begin{description}

  \item[abstract]
    Abstract of the item (e.g. the abstract of a journal article)

  \item[annote]
    Short markup, decoration, or annotation to the item (e.g., to indicate
    items included in a review);\\
    For descriptive text (e.g., in an annotated bibliography), use
    \texttt{note} instead

  \item[archive]
    Archive storing the item

  \item[archive\_collection]
    Collection the item is part of within an archive

  \item[archive\_location]
    Storage location within an archive (e.g. a box and folder number)

  \item[archive-place]
    Geographic location of the archive

  \item[authority]
    Issuing or judicial authority (e.g. "USPTO" for a patent, "Fairfax
    Circuit Court" for a legal case)

  \item[call-number]
    Call number (to locate the item in a library)

  \item[citation-key]
    Identifier of the item in the input data file (analogous to BibTeX
    entrykey);\\
    Use this variable to facilitate conversion between word-processor and
    plain-text writing systems;\\
    For an identifer intended as formatted output label for a citation (e.g.
    "Ferr78"), use \texttt{citation-label} instead

  \item[citation-label]
    Label identifying the item in in-text citations of label styles (e.g.
    "Ferr78");\\
    May be assigned by the CSL processor based on item metadata;\\
    For the identifier of the item in the input data file, use
    \texttt{citation-key} instead

  \item[collection-title]
    Title of the collection holding the item (e.g. the series title for a
    book; the lecture series title for a presentation)

  \item[container-title]
    Title of the container holding the item (e.g. the book title for a book
    chapter, the journal title for a journal article; the album title for a
    recording; the session title for multi-part presentation at a
    conference)

  \item[container-title-short]
    Short/abbreviated form of \texttt{container-title};\\
    Deprecated; use \texttt{variable="container-title"\ form="short"}
    instead

  \item[dimensions]
    Physical (e.g. size) or temporal (e.g. running time) dimensions of the
    item

  \item[division]
    Minor subdivision of a court with a \texttt{jurisdiction} for a legal
    item

  \item[DOI]
    Digital Object Identifier (e.g. "10.1128/AEM.02591-07")

  \item[event]
    Deprecated legacy variant of \texttt{event-title}

  \item[event-title]
    Name of the event related to the item (e.g. the conference name when
    citing a conference paper; the meeting where presentation was made)

  \item[event-place]
    Geographic location of the event related to the item (e.g. "Amsterdam,
    The Netherlands")

  \item[genre]
    Type, class, or subtype of the item (e.g. "Doctoral dissertation" for a
    PhD thesis; "NIH Publication" for an NIH technical report);\\
    Do not use for topical descriptions or categories (e.g. "adventure" for
    an adventure movie)

  \item[ISBN]
    International Standard Book Number (e.g. "978-3-8474-1017-1")

  \item[ISSN]
    International Standard Serial Number

  \item[jurisdiction]
    Geographic scope of relevance (e.g. "US" for a US patent; the court
    hearing a legal case)

  \item[keyword]
    Keyword(s) or tag(s) attached to the item

  \item[language]
    The language of the item;\\
    Should be entered as an ISO 639-1 two-letter language code (e.g. "en",
    "zh"), optionally with a two-letter locale code (e.g. "de-DE", "de-AT")

  \item[license]
    The license information applicable to an item (e.g. the license an
    article or software is released under; the copyright information for an
    item; the classification status of a document)

  \item[medium]
    Description of the item\textquotesingle s format or medium (e.g. "CD",
    "DVD", "Album", etc.)

  \item[note]
    Descriptive text or notes about an item (e.g. in an annotated
    bibliography)

  \item[original-publisher]
    Original publisher, for items that have been republished by a different
    publisher

  \item[original-publisher-place]
    Geographic location of the original publisher (e.g. "London, UK")

  \item[original-title]
    Title of the original version (e.g. "Война и мир", the untranslated
    Russian title of "War and Peace")

  \item[part-title]
    Title of the specific part of an item being cited

  \item[PMCID]
    PubMed Central reference number

  \item[PMID]
    PubMed reference number

  \item[publisher]
    Publisher

  \item[publisher-place]
    Geographic location of the publisher

  \item[references]
    Resources related to the procedural history of a legal case or
    legislation;\\
    Can also be used to refer to the procedural history of other items (e.g.
    "Conference canceled" for a presentation accepted as a conference that
    was subsequently canceled; details of a retraction or correction notice)

  \item[reviewed-genre]
    Type of the item being reviewed by the current item (e.g. book, film)

  \item[reviewed-title]
    Title of the item reviewed by the current item

  \item[scale]
    Scale of e.g. a map or model

  \item[source]
    Source from whence the item originates (e.g. a library catalog or
    database)

  \item[status]
    Publication status of the item (e.g. "forthcoming"; "in press"; "advance
    online publication"; "retracted")

  \item[title]
    Primary title of the item

  \item[title-short]
    Short/abbreviated form of \texttt{title};\\
    Deprecated; use \texttt{variable="title"\ form="short"} instead

  \item[URL]
    Uniform Resource Locator (e.g.
    "\url{https://aem.asm.org/cgi/content/full/74/9/2766}")

  \item[volume-title]
    Title of the volume of the item or container holding the item;\\
    Also use for titles of periodical special issues, special sections, and
    the like

  \item[year-suffix]
    Disambiguating year suffix in author-date styles (e.g. "a" in "Doe,
    1999a")
    \end{description}

\hypertarget{number-variables}{%
\paragraph{Number Variables}\label{number-variables}}

Number variables are a subset of the
\protect\hyperlink{standard-variables}{Standard Variables}.

\begin{description}

  \item[chapter-number]
    Chapter number (e.g. chapter number in a book; track number on an album)

  \item[citation-number]
    Index (starting at 1) of the cited reference in the bibliography
    (generated by the CSL processor)

  \item[collection-number]
    Number identifying the collection holding the item (e.g. the series
    number for a book)

  \item[edition]
    (Container) edition holding the item (e.g. "3" when citing a chapter in
    the third edition of a book)

  \item[first-reference-note-number]
    Number of a preceding note containing the first reference to the item;\\
    Assigned by the CSL processor;\\
    Empty in non-note-based styles or when the item hasn\textquotesingle t
    been cited in any preceding notes in a document

  \item[issue]
    Issue number of the item or container holding the item (e.g. "5" when
    citing a journal article from journal volume 2, issue 5);\\
    Use \texttt{volume-title} for the title of the issue, if any

  \item[locator]
    A cite-specific pinpointer within the item (e.g. a page number within a
    book, or a volume in a multi-volume work);\\
    Must be accompanied in the input data by a label indicating the locator
    type (see the \protect\hyperlink{locators}{Locators} term list), which
    determines which term is rendered by \texttt{cs:label} when the
    \texttt{locator} variable is selected.

  \item[number]
    Number identifying the item (e.g. a report number)

  \item[number-of-pages]
    Total number of pages of the cited item

  \item[number-of-volumes]
    Total number of volumes, used when citing multi-volume books and such

  \item[page]
    Range of pages the item (e.g. a journal article) covers in a container
    (e.g. a journal issue)

  \item[page-first]
    First page of the range of pages the item (e.g. a journal article)
    covers in a container (e.g. a journal issue)

  \item[part-number]
    Number of the specific part of the item being cited (e.g. part 2 of a
    journal article);\\
    Use \texttt{part-title} for the title of the part, if any

  \item[printing-number]
    Printing number of the item or container holding the item

  \item[section]
    Section of the item or container holding the item (e.g. "§2.0.1" for a
    law; "politics" for a newspaper article)

  \item[supplement-number]
    Supplement number of the item or container holding the item (e.g. for
    secondary legal items that are regularly updated between editions)

  \item[version]
    Version of the item (e.g. "2.0.9" for a software program)

  \item[volume]
    Volume number of the item (e.g. "2" when citing volume 2 of a book) or
    the container holding the item (e.g. "2" when citing a chapter from
    volume 2 of a book);\\
    Use \texttt{volume-title} for the title of the volume, if any
    \end{description}

\hypertarget{date-variables}{%
\subsubsection{Date Variables}\label{date-variables}}

\begin{description}

  \item[accessed]
    Date the item has been accessed

  \item[available-date]
    Date the item was initially available (e.g. the online publication date
    of a journal article before its formal publication date; the date a
    treaty was made available for signing)

  \item[event-date]
    Date the event related to an item took place

  \item[issued]
    Date the item was issued/published

  \item[original-date]
    Issue date of the original version

  \item[submitted]
    Date the item (e.g. a manuscript) was submitted for publication
    \end{description}

\hypertarget{name-variables}{%
\subsubsection{Name Variables}\label{name-variables}}

\begin{description}

  \item[author]
    Author

  \item[chair]
    The person leading the session containing a presentation (e.g. the
    organizer of the \texttt{container-title} of a \texttt{speech})

  \item[collection-editor]
    Editor of the collection holding the item (e.g. the series editor for a
    book)

  \item[compiler]
    Person compiling or selecting material for an item from the works of
    various persons or bodies (e.g. for an anthology)

  \item[composer]
    Composer (e.g. of a musical score)

  \item[container-author]
    Author of the container holding the item (e.g. the book author for a
    book chapter)

  \item[contributor]
    A minor contributor to the item; typically cited using "with" before the
    name when listed in a bibliography

  \item[curator]
    Curator of an exhibit or collection (e.g. in a museum)

  \item[director]
    Director (e.g. of a film)

  \item[editor]
    Editor

  \item[editorial-director]
    Managing editor ("Directeur de la Publication" in French)

  \item[editor-translator]
    Combined editor and translator of a work;\\
    The citation processory must be automatically generate if
    \texttt{editor} and \texttt{translator} variables are identical;\\
    May also be provided directly in item data

  \item[executive-producer]
    Executive producer (e.g. of a television series)

  \item[guest]
    Guest (e.g. on a TV show or podcast)

  \item[host]
    Host (e.g. of a TV show or podcast)

  \item[illustrator]
    Illustrator (e.g. of a children\textquotesingle s book or graphic novel)

  \item[interviewer]
    Interviewer (e.g. of an interview)

  \item[narrator]
    Narrator (e.g. of an audio book)

  \item[organizer]
    Organizer of an event (e.g. organizer of a workshop or conference)

  \item[original-author]
    The original creator of a work (e.g. the form of the author name listed
    on the original version of a book; the historical author of a work; the
    original songwriter or performer for a musical piece; the original
    developer or programmer for a piece of software; the original author of
    an adapted work such as a book adapted into a screenplay)

  \item[performer]
    Performer of an item (e.g. an actor appearing in a film; a muscian
    performing a piece of music)

  \item[producer]
    Producer (e.g. of a television or radio broadcast)

  \item[recipient]
    Recipient (e.g. of a letter)

  \item[reviewed-author]
    Author of the item reviewed by the current item

  \item[script-writer]
    Writer of a script or screenplay (e.g. of a film)

  \item[series-creator]
    Creator of a series (e.g. of a television series)

  \item[translator]
    Translator
    \end{description}

\hypertarget{appendix-v---page-range-formats}{%
\subsection{Appendix V - Page Range
Formats}\label{appendix-v---page-range-formats}}

The page abbreviation rules for the different values of the
\texttt{page-range-format} attribute on \texttt{cs:style} are:

\begin{description}
\item["chicago"]
Alias for "chicago-15"; will change to be an alias for "chicago-16" in
CSL v1.1.
\item["chicago-15"]
Page ranges are abbreviated according to the \href{}{Chicago Manual of
Style (15th ed and earlier) rules} (see 15th ed, section 9.64):

\begin{longtable}[]{@{}
  >{\raggedright\arraybackslash}p{(\columnwidth - 4\tabcolsep) * \real{0.3125}}
  >{\raggedright\arraybackslash}p{(\columnwidth - 4\tabcolsep) * \real{0.3375}}
  >{\raggedright\arraybackslash}p{(\columnwidth - 4\tabcolsep) * \real{0.3375}}@{}}
\toprule\noalign{}
\begin{minipage}[b]{\linewidth}\raggedright
First number
\end{minipage} & \begin{minipage}[b]{\linewidth}\raggedright
Second number
\end{minipage} & \begin{minipage}[b]{\linewidth}\raggedright
Examples
\end{minipage} \\
\midrule\noalign{}
\endhead
\bottomrule\noalign{}
\endlastfoot
Less than 100 & Use all digits & 3 -- 10; 71 -- 72 \\
100 or multiple of 100 & Use all digits & 100 -- 104; 600 -- 613; 1100
-- 1123 \\
101 through 109 (in multiples of 100) & Use changed part only, omitting
unneeded zeros & 107 -- 8; 505 -- 17; 1002 -- 6 \\
110 through 199 (in multiples of 100) & Use two digits, or more as
needed & 321 -- 25; 415 -- 532; 11564 -- 68; 13792 -- 803 \\
4 digits & If numbers are four digits long and three digits change, use
all digits & 1496 -- 1504; 2787 -- 2816 \\
\end{longtable}
\item["chicago-16"]
Page ranges are abbreviated according to the
\href{http://cmosshoptalk.com/2018/04/10/316-7-316-17-or-316-317-chicago-style-for-number-ranges/}{Chicago
Manual of Style (16th ed and later) rules} (see 16th ed, section 9.61):

\begin{longtable}[]{@{}
  >{\raggedright\arraybackslash}p{(\columnwidth - 4\tabcolsep) * \real{0.3947}}
  >{\raggedright\arraybackslash}p{(\columnwidth - 4\tabcolsep) * \real{0.3553}}
  >{\raggedright\arraybackslash}p{(\columnwidth - 4\tabcolsep) * \real{0.2368}}@{}}
\toprule\noalign{}
\begin{minipage}[b]{\linewidth}\raggedright
First number
\end{minipage} & \begin{minipage}[b]{\linewidth}\raggedright
Second number
\end{minipage} & \begin{minipage}[b]{\linewidth}\raggedright
Examples
\end{minipage} \\
\midrule\noalign{}
\endhead
\bottomrule\noalign{}
\endlastfoot
Less than 100 & Use all digits & 3 -- 10; 71 -- 72; 92 -- 113; \\
100 or multiple of 100 & Use all digits & 100 -- 104; 600 -- 613; 1100
-- 1123 \\
101 through 109, 201 through 209, etc. (for each multiple of 100) & Use
changed part only, omitting unneeded zeros & 107 -- 8; 505 -- 17; 1002
-- 6 \\
Everything else (110 through 199, 210 through 299, etc.; for each
multiple of 100) & Use two digits, unless more digits are needed to show
the changed part & 321 -- 25; 415 -- 532; 1087 -- 89; 1496 -- 500; 11564
-- 68; 13792 -- 803 12991 -- 3001 \\
\end{longtable}
\item["expanded"]
Abbreviated page ranges are expanded to their non-abbreviated form: 42
-- 45, 321 -- 328, 2787 -- 2816
\item["minimal"]
All digits repeated in the second number are left out: 42 -- 5, 321 --
8, 2787 -- 816
\item["minimal-two"]
As "minimal", but at least two digits are kept in the second number when
it has two or more digits long.
\end{description}

\hypertarget{appendix-vi-links}{%
\subsection{Appendix VI: Links}\label{appendix-vi-links}}

The CSL syntax does not have support for configuration of links.
However, processors should include links on bibliographic references,
using the following rules:

If the bibliography entry for an item renders any of the following
identifiers, the identifier should be anchored as a link, with the
target of the link as follows:

\begin{enumerate}
\def\labelenumi{\arabic{enumi}.}
\tightlist
\item
  \texttt{url}: output as is
\item
  \texttt{doi}: prepend with "\url{https://doi.org/}"
\item
  \texttt{pmid}: prepend with
  "\url{https://www.ncbi.nlm.nih.gov/pubmed/}"
\item
  \texttt{pmcid}: prepend with
  "\url{https://www.ncbi.nlm.nih.gov/pmc/articles/}"
\end{enumerate}

If the identifier is rendered as a URI, include rendered URI components
(e.g. "\url{https://doi.org/}") in the link anchor. Do not include any
other affix text in the link anchor (e.g. "Available from: ", "doi: ",
"PMID: ").

Citation processors should include an option flag for calling
applications to disable bibliography linking behavior.
